% Luc Destombe --- licence CC BY-NC-SA 4.0
% https://creativecommons.org/licenses/by-nc-sa/4.0/deed.fr
% Fichier autonome : compiler avec pdflatex, deux fois (signets, nombre de pages).
\documentclass[11pt,a4paper]{article}
\makeatletter
\newif\iflycee@unecolonne
\lycee@unecolonnetrue
\newif\iflycee@palatino
\lycee@palatinotrue

\RequirePackage[utf8]{inputenc}
\RequirePackage[T1]{fontenc}
\RequirePackage[french]{babel}
\RequirePackage{etoolbox}

\newcommand{\ShorthandsOff}{\@ifundefined{shorthandoff}{}{\shorthandoff{;:!?}}}
\newcommand{\ShorthandsOn}{\@ifundefined{shorthandon}{}{\shorthandon{;:!?}}}
\ShorthandsOff
\AtBeginDocument{\ShorthandsOn}
\AtBeginEnvironment{tikzpicture}{\ShorthandsOff}

\RequirePackage{geometry}
\RequirePackage{fancyhdr}
\RequirePackage{lastpage}
\RequirePackage{multicol}
\RequirePackage{array}
\RequirePackage{enumitem}
\RequirePackage{xcolor}
\RequirePackage{needspace}

\RequirePackage{amsmath,amssymb}
\RequirePackage{mathpazo}
\DeclareMathSizes{10.95}{10.95}{8}{5.5}
\begingroup\catcode`\_=\active \catcode`\^=\active
\gdef_#1{\sb{\scriptscriptstyle #1}}
\gdef^#1{\sp{\scriptscriptstyle\lycee@moinsepais #1}}
\endgroup
\newcommand\lycee@moins{\mathbin{%
  \setbox\z@\hbox{$\m@th\scriptscriptstyle\lycee@moinsorig$}%
  \dimen@=.55\wd\z@ \advance\dimen@-.5pt
  \hbox to.75\wd\z@{\hss$\m@th\scriptscriptstyle\vcenter{\hbox{%
    \tikz\draw[line width=.5pt,line cap=round](0,0)--(\the\dimen@,0);}}$\hss}}}
\begingroup\lccode`\~=`\- \lowercase{\endgroup\let~}\lycee@moins
\newcommand\lycee@moinsepais{\mathcode`\-="8000 }
\AtBeginDocument{\mathchardef\lycee@moinsorig=\mathcode`\-\relax}
\AtBeginDocument{\catcode`\_=12 \mathcode`\_="8000
  \catcode`\^=12 \mathcode`\^="8000 }
\newcommand\lycee@exposants{%
  \fontdimen13\textfont2=.48\fontdimen6\textfont2
  \fontdimen14\textfont2=.48\fontdimen6\textfont2
  \fontdimen15\textfont2=.48\fontdimen6\textfont2
  \fontdimen13\scriptfont2=.48\fontdimen6\scriptfont2
  \fontdimen14\scriptfont2=.48\fontdimen6\scriptfont2
  \fontdimen15\scriptfont2=.48\fontdimen6\scriptfont2}
\every@math@size\expandafter{\the\every@math@size\lycee@exposants}
\newlength{\retraitcalculs}
\setlength{\retraitcalculs}{2em}

\RequirePackage{tikz}
\usetikzlibrary{arrows.meta,calc,babel,shapes.misc}
\RequirePackage[most]{tcolorbox}
\tcbuselibrary{skins,breakable}

\definecolor{coulDef}{HTML}{1F4E79}
\definecolor{coulProp}{HTML}{7030A0}
\definecolor{coulRem}{HTML}{C55A11}
\definecolor{coulEx}{HTML}{2E7D32}
\definecolor{coulExo}{HTML}{B03A2E}
\definecolor{coulPart}{HTML}{1F4E79}
\definecolor{coulMeth}{HTML}{1B6E4F}
\definecolor{coulCourbe}{HTML}{0B5ED7}
\definecolor{coulCourbeB}{HTML}{C00000}
\definecolor{coulCourbeC}{HTML}{2E7D32}
\definecolor{coulGrille}{HTML}{8C8C8C}
\definecolor{coulRep}{HTML}{1B6E4F}
\definecolor{coulBar}{HTML}{2E7D32}

\newcommand{\R}{\mathbb{R}}
\newcommand{\Rs}{\mathbb{R}^{*}}
\renewcommand{\le}{\leqslant}
\renewcommand{\ge}{\geqslant}
\renewcommand{\approx}{\simeq}
\newcommand{\vect}[1]{\overrightarrow{#1}}
\newcommand{\norme}[1]{\left\lVert #1 \right\rVert}
\newcommand{\intervalle}[2]{\left[#1\,;#2\right]}
\RequirePackage{cancel}
\renewcommand{\CancelColor}{\color{coulExo}}
\newcommand{\fois}[1]{\times{\color{coulExo}#1}}
\newcommand{\barre}[1]{\cancel{#1}}

\tikzset{grille/.style={coulGrille,line width=0.4pt}}
\newcommand{\quadrillage}[4]{\draw[grille,step=1] (#1,#2) grid (#3,#4);}
\tikzset{
  croix/.style={cross out,draw,line width=0.6pt,inner sep=0pt,minimum size=4pt},
  croixdroite/.style={croix,rotate=45,line width=0.8pt},
}
\newcommand{\pt}[3]{\path (#1) node[croix]{} node[#2,font=\small,inner sep=2.5pt]{$#3$};}

\newcommand{\lycee@repere}[4]{%
  \draw[grille,step=1] (#1,#2) grid (#3,#4);
  \draw[-{Stealth[length=2mm]},thick] (#1,0)--({#3+0.4},0) node[below left=-1pt]{$x$};
  \draw[-{Stealth[length=2mm]},thick] (0,#2)--(0,{#4+0.4}) node[below left=-1pt]{$y$};
  \node[below left,font=\scriptsize,inner sep=1.5pt] at (0,0) {$0$};}
\newcommand{\lycee@gradx}[1]{\foreach \k/\l in {#1}{%
  \draw (\k,0.07)--(\k,-0.07) node[below,font=\scriptsize,inner sep=1.5pt]{$\l$};}}
\newcommand{\lycee@grady}[1]{\foreach \k/\l in {#1}{%
  \draw (0.07,\k)--(-0.07,\k) node[left,font=\scriptsize,inner sep=1.5pt]{$\l$};}}
\newcommand{\lycee@intff}[2]{\left[#1\,;#2\right]}
\newcommand{\lycee@intfo}[2]{\left[#1\,;#2\right[}
\newcommand{\lycee@intof}[2]{\left]#1\,;#2\right]}
\newcommand{\lycee@intoo}[2]{\left]#1\,;#2\right[}
\newcommand{\lycee@ens}[1]{\left\{#1\right\}}
\AtBeginDocument{%
  \providecommand{\repere}{\lycee@repere}%
  \providecommand{\gradx}{\lycee@gradx}%
  \providecommand{\grady}{\lycee@grady}%
  \providecommand{\Cf}{\mathcal{C}\sb{\scriptscriptstyle f}}%
  \providecommand{\Cg}{\mathcal{C}\sb{\scriptscriptstyle g}}%
  \providecommand{\intff}{\lycee@intff}%
  \providecommand{\intfo}{\lycee@intfo}%
  \providecommand{\intof}{\lycee@intof}%
  \providecommand{\intoo}{\lycee@intoo}%
  \providecommand{\ens}{\lycee@ens}}

\newlist{questions}{enumerate}{3}
\setlist[questions,1]{label=\arabic*\textdegree),leftmargin=*,itemsep=2pt,topsep=3pt}
\setlist[questions,2]{label=\alph*),leftmargin=*,itemsep=1pt,topsep=2pt}
\setlist[questions,3]{label=\roman*),leftmargin=*,itemsep=1pt,topsep=1pt}
\setlist[itemize]{label=\textbullet,leftmargin=*,itemsep=2pt,topsep=3pt}

\newcommand{\besoinplace}[1]{\par\penalty\@M\begingroup
  \dimen@=#1\relax\dimen@ii=\pagegoal\advance\dimen@ii-\pagetotal
  \ifdim\dimen@>\dimen@ii\ifdim\dimen@ii>\z@\vfil\fi\break\fi\endgroup}

\newcommand{\lycee@pied}{\thepage/\pageref{LastPage}}
\newcommand{\entete}[2]{%
  \pagestyle{fancy}\fancyhf{}%
  \fancyhead[L]{\footnotesize\color{coulPart}\textbf{#1}}%
  \fancyhead[R]{\footnotesize\color{coulPart}\textbf{#2}}%
  \fancyfoot[C]{\footnotesize\lycee@pied}%
  \renewcommand{\headrulewidth}{0.4pt}%
  \renewcommand{\headrule}{\hbox to\headwidth{%
    \color{coulPart!40}\leaders\hrule height \headrulewidth\hfill}}}

\RequirePackage{ccicons}
\newcommand{\lycee@licence}{{\scriptsize\color{gray}%
  \copyright\ Luc Destombe --- \ccbyncsa\ CC BY-NC-SA 4.0}}
\newif\iflycee@licence \lycee@licencetrue
\newcommand{\sanslicence}{\lycee@licencefalse}
\AtBeginDocument{\iflycee@licence\AddToHookNext{shipout/foreground}{%
  \put(0,0){\rlap{\kern\dimexpr1in+\hoffset+\oddsidemargin+\textwidth\relax
    \llap{\raisebox{-\dimexpr1in+\voffset+\topmargin+\headheight+\headsep
      +\textheight+\footskip\relax}{\lycee@licence}}}}}\fi}

\newcommand{\formulesserrees}{%
  \setlength{\abovedisplayskip}{3pt}\setlength{\belowdisplayskip}{3pt}%
  \setlength{\abovedisplayshortskip}{2pt}\setlength{\belowdisplayshortskip}{3pt}}

\setlength{\parindent}{0pt}

\RequirePackage[implicit=false,hidelinks,pdfencoding=auto,psdextra,bookmarksopen,
  bookmarksopenlevel=1,pdfcreator={},pdfproducer={}]{hyperref}
\RequirePackage{bookmark}
\hypersetup{pdfauthor={Luc Destombe},pdfkeywords={CC BY-NC-SA 4.0}}
\newcounter{lycee@signet}
\newcommand{\signet}[3][]{\stepcounter{lycee@signet}%
  \edef\lycee@dest{\if\relax\detokenize{#1}\relax
    signet.\arabic{lycee@signet}\else#1\fi}%
  \hypertarget{\lycee@dest}{}\bookmark[level=#2,dest=\lycee@dest]{#3}}
\pdfstringdefDisableCommands{%
  \def\R{R}\def\vect#1{#1}\def\Cf{Cf}\def\Cg{Cg}\def\fois#1{×#1}%
  \def\barre#1{#1}\def\intervalle#1#2{[#1;#2]}\def\intff#1#2{[#1;#2]}%
  \def\textsuperscript#1{#1}\def\dfrac#1#2{#1/#2}\def\frac#1#2{#1/#2}%
  \def\vec#1{#1⃗}}%
\RequirePackage[official]{eurosym}

\tcbset{exercice corrige/.style={
  enhanced,breakable,before upper=\raggedright,
  colback=white,colframe=coulExo,
  boxrule=0.8pt,arc=2pt,
  left=6pt,right=6pt,top=8pt,bottom=5pt,
  before skip=9pt,after skip=4pt,
  coltitle=white,fonttitle=\bfseries\small,
  attach boxed title to top left={xshift=8pt,yshift=-\tcboxedtitleheight/2},
  boxed title style={colback=coulExo,arc=2pt,boxrule=0pt}}}
\newcounter{exo}
\newtcolorbox{exercice}[1][]{exercice corrige,
  phantom={\signet[exercice-\arabic{exo}]{2}{Exercice \arabic{exo}}},
  title={Exercice \theexo\ifx\relax#1\relax\else{} --- #1\fi}}
\newenvironment{exo}[1][\relax]{\refstepcounter{exo}\begin{exercice}[#1]}{\end{exercice}}

\geometry{top=1.8cm,bottom=1cm,left=1cm,right=1cm,headheight=14pt,footskip=0.6cm}
\renewcommand{\lycee@pied}{\thepage}

\renewcommand{\theexo}{\ifnum\value{exo}<10 0\fi\arabic{exo}}
\tcbset{exercice corrige/.style={
  enhanced,breakable,before upper=\raggedright,
  colback=white,colframe=coulExo,
  boxrule=0.9pt,arc=2pt,
  left=8pt,right=8pt,top=8pt,bottom=6pt,
  before skip=8pt,after skip=6pt,
  coltitle=white,fonttitle=\bfseries,
  attach boxed title to top left={xshift=8pt,yshift=-\tcboxedtitleheight/2},
  boxed title style={colback=coulExo,arc=2pt,boxrule=0pt}}}

\newcommand{\entetefiche}[2]{%
  \begin{center}
    \begin{tcolorbox}[enhanced,width=0.92\linewidth,colback=white,colframe=coulPart,
        boxrule=1.4pt,arc=3pt,halign=center,top=6pt,bottom=6pt]
      {\large\bfseries\color{coulPart} CORRIGÉ --- FICHE D'EXERCICES}\\[3pt]
      {\normalsize\bfseries #1}\\[2pt]
      {\normalsize #2}
    \end{tcolorbox}
  \end{center}}

\newcounter{partie}
\newcommand{\partie}[1]{%
  \stepcounter{partie}%
  \besoinplace{8\baselineskip}\vspace{4pt}%
  \begin{tcolorbox}[enhanced,colback=coulPart!8,colframe=coulPart,boxrule=0pt,
      leftrule=3pt,arc=0pt,left=6pt,right=4pt,top=3pt,bottom=3pt,
      before skip=6pt,after skip=2pt]
    \bfseries\color{coulPart}\Roman{partie} --- #1\signet{1}{\Roman{partie} --- #1}
  \end{tcolorbox}}
\newcommand{\partiesansnum}[1]{%
  \besoinplace{8\baselineskip}\vspace{4pt}%
  \begin{tcolorbox}[enhanced,colback=coulPart!8,colframe=coulPart,boxrule=0pt,
      leftrule=3pt,arc=0pt,left=6pt,right=4pt,top=3pt,bottom=3pt,
      before skip=6pt,after skip=2pt]
    \bfseries\color{coulPart}#1\signet{1}{#1}
  \end{tcolorbox}}

\setlist[questions,1]{label=\arabic*\textdegree),leftmargin=*,itemsep=3pt,topsep=3pt}
\setlist[questions,2]{label=\alph*),leftmargin=*,itemsep=2pt,topsep=2pt}
\setlist[questions,3]{label=\roman*),leftmargin=*,itemsep=1pt,topsep=1pt}
\newlist{expr}{enumerate}{1}
\setlist[expr]{label={\Alph*\ $=$},leftmargin=*,itemsep=2pt,topsep=3pt}

\newcommand{\res}[1]{{\color{coulRep}#1}}
\newcommand{\rem}[1]{\par\smallskip{\small\itshape\color{black!60}$\rightarrow$ #1}}
\newenvironment{calculs}{\par\smallskip\noindent$\displaystyle\begin{aligned}[t]}%
  {\end{aligned}$\par\smallskip}
\newcommand{\justif}[1]{\quad\text{\small\itshape\color{black!60}#1}}
\newcommand{\cas}[1]{\par\smallskip\noindent\textbf{\color{coulExo}#1}\nobreak}
\newcommand{\explique}[2]{\par\smallskip\noindent
  \begin{minipage}[t]{0.6\linewidth}\raggedright #1\end{minipage}\hfill
  \begin{minipage}[t]{0.37\linewidth}\raggedright\leavevmode
    {\small\itshape\color{black!60}#2}\end{minipage}\par}

\setlength{\parskip}{2pt}

\makeatother
\entete{Chapitre 5 --- Suites numériques --- corrigé}{Première spécialité}
\usepackage{tkz-tab}

\let\casligne\cas
\renewcommand{\cas}[1]{\casligne{#1}\par\nopagebreak}
\setlength{\parindent}{0pt}

\newtcblisting{python}{%
  listing only,enhanced,
  colback=black!3,colframe=black!35,boxrule=0.5pt,arc=1pt,
  left=4pt,right=4pt,top=1pt,bottom=1pt,
  listing options={language=Python,basicstyle=\ttfamily\small,
    keywordstyle=\color{coulMeth}\bfseries,columns=fullflexible,
    keepspaces=true,showstringspaces=false,aboveskip=0pt,belowskip=0pt}}
\newcommand{\reperesuite}[4]{%
  \draw[grille,step=1] (-1,#2) grid (#1,#3);
  \draw[-{Stealth[length=2mm]},thick] (-1,0)--({#1+0.4},0)
    node[below left=-1pt,font=\footnotesize]{$n$};
  \draw[-{Stealth[length=2mm]},thick] (0,#2)--(0,{#3+0.4})
    node[below left=-1pt,font=\footnotesize]{$#4$};
  \node[below left,font=\scriptsize,inner sep=1.5pt] at (0,0) {$0$};}
\newcommand{\gris}[1]{{\small\itshape\color{black!60}#1}}

\begin{document}

\entetefiche{Chapitre 5 --- Suites numériques}{Première spécialité mathématiques}

\partie{Définition d'une suite}

\begin{exo}[Repérer les termes et les rangs]
\explique{Le premier terme est $u_{0}$ : on numérote les termes à partir de $0$.}{le
rang est écrit au-dessus de chaque terme}
\begin{center}
\renewcommand{\arraystretch}{1.3}
\begin{tabular}{|c|c|c|c|c|c|c|}
\hline
rang $n$ & $0$ & $1$ & $2$ & $3$ & $4$ & $5$\\\hline
terme $u_{n}$ & $5$ & $9$ & $13$ & $17$ & $21$ & $25$\\\hline
\end{tabular}
\end{center}
\cas{1.}
\explique{$\res{u_{0}=5}$ et $\res{u_{3}=17}$.}{on lit sous $n=0$ et sous $n=3$}
\cas{2.}
\explique{$21=u_{4}$ : le terme $21$ est de \res{rang $4$}.}{le rang est
l'indice, pas la valeur du terme}
\cas{3.}
\explique{Le 4\textsuperscript{e} terme est $\res{u_{3}=17}$.}{on compte $u_{0}$,
$u_{1}$, $u_{2}$ puis $u_{3}$}
\rem{Erreur fréquente : répondre $u_{4}$. Quand la suite commence à $u_{0}$, le
4\textsuperscript{e} terme porte l'indice $3$ : l'indice a toujours un de moins que la
place.}
\cas{4.}
\explique{Le terme qui suit $u_{3}$ est $\res{u_{4}=21}$ ; celui qui le précède est
$\res{u_{2}=13}$.}{$u_{3+1}$ et $u_{3-1}$ : on avance ou on recule d'un rang}
\end{exo}

\begin{exo}[Une suite qui commence à $u_{1}$]
\explique{Le premier terme est $v_{1}$ : on numérote à partir de $1$.}{il n'y a pas de
$v_{0}$}
\begin{center}
\renewcommand{\arraystretch}{1.3}
\begin{tabular}{|c|c|c|c|c|c|}
\hline
rang $n$ & $1$ & $2$ & $3$ & $4$ & $5$\\\hline
terme $v_{n}$ & $1$ & $4$ & $9$ & $16$ & $25$\\\hline
\end{tabular}
\end{center}
\cas{1.}
\explique{$\res{v_{1}=1}$ et $\res{v_{4}=16}$.}{on lit sous $n=1$ et sous $n=4$}
\cas{2.}
\explique{Le 5\textsuperscript{e} terme est $\res{v_{5}=25}$ : il est de
\res{rang $5$}.}{ici, la place et le rang coïncident}
\rem{À comparer avec l'exercice 1 : tout dépend du premier rang. Avec $u_{0}$, le
5\textsuperscript{e} terme est $u_{4}$ ; avec $v_{1}$, c'est $v_{5}$.}
\cas{3.}
\explique{On compare chaque terme à son rang : $v_{1}=1=1^{2}$ ; $v_{2}=4=2^{2}$ ;
$v_{3}=9=3^{2}$ ; $v_{4}=16=4^{2}$ ; $v_{5}=25=5^{2}$.}{chaque terme est le carré de
son rang}
\explique{On conjecture que, pour tout $n\ge 1$, $\res{v_{n}=n^{2}}$.}{une conjecture :
on ne connaît que cinq termes}
\end{exo}

\begin{exo}[Premier rang]
\explique{Pour chaque suite, on cherche les rangs $n$ pour lesquels le calcul de $u_{n}$
est possible, puis le plus petit d'entre eux.}{deux interdits : diviser par $0$,
prendre la racine d'un nombre négatif}
\cas{a)}
\explique{On ne divise pas par $0$ : il faut $n\neq 0$.}{valeur interdite}
\explique{$(u_{n})$ est définie à partir du rang $\res{n_{0}=1}$.}{premier entier
différent de $0$}
\cas{b)}
\explique{Sous la racine, il faut $n-3\ge 0$, soit $n\ge 3$.}{une racine carrée
n'existe que pour un nombre positif ou nul}
\explique{$(v_{n})$ est définie à partir du rang $\res{n_{0}=3}$.}{$v_{3}=\sqrt{0}=0$
existe bien}
\cas{c)}
\explique{Il faut $n-5\ge 0$ (racine) et $\sqrt{n-5}\neq 0$ (division), soit
$n-5>0$.}{les deux conditions à la fois}
\explique{$n>5$ : $(w_{n})$ est définie à partir du rang $\res{n_{0}=6}$.}{$n$ est
entier : après $5$ vient $6$}
\rem{Erreur fréquente : répondre $n_{0}=5$. Mais $w_{5}=\dfrac{4}{\sqrt{0}}=\dfrac{4}{0}$
n'existe pas.}
\cas{d)}
\explique{Racine : il faut $n\ge 0$, ce qui est vrai pour tout entier naturel.
Division : il faut $n\neq 0$.}{on regarde chaque morceau de l'expression}
\explique{$(t_{n})$ est définie à partir du rang $\res{n_{0}=1}$.}{$t_{1}=\sqrt{1}+1=2$}
\end{exo}

\begin{exo}[Les décimales de $\frac{1}{7}$]
\cas{1.}
\explique{On numérote les chiffres après la virgule de $0{,}142857142857\dots$ :}{le
$n$-ième chiffre est $u_{n}$}
\begin{center}
\renewcommand{\arraystretch}{1.3}
\begin{tabular}{|c|c|c|c|c|c|c|c|c|}
\hline
$n$ & $1$ & $2$ & $3$ & $4$ & $5$ & $6$ & $7$ & $8$\\\hline
$u_{n}$ & $1$ & $4$ & $\res{2}$ & $\res{8}$ & $\res{5}$ & $\res{7}$ & $\res{1}$ &
$\res{4}$\\\hline
\end{tabular}
\end{center}
\explique{$u_{3}=2$ ; $u_{4}=8$ ; $u_{5}=5$ ; $u_{6}=7$ ; $u_{7}=1$ ;
$u_{8}=4$.}{lecture des décimales}
\explique{Le premier rang de la suite est $\res{1}$.}{$u_{1}$ est la première décimale ;
il n'y a pas de $u_{0}$}
\cas{2.}
\explique{Le motif $1$ ; $4$ ; $2$ ; $8$ ; $5$ ; $7$ se répète : $u_{7}=u_{1}$,
$u_{8}=u_{2}$\dots{} Les termes reviennent tous les $6$ rangs.}{la période est $6$}
\explique{On compte combien de motifs complets tiennent avant le rang cherché : c'est
la division euclidienne par $6$.}{le reste donne la place dans le motif}
\begin{calculs}
20 &= 3\times 6+2 &&\justif{trois motifs complets, puis $2$ rangs}\\
u_{20} &= u_{2} = \res{4} &&\justif{la 2\textsuperscript{e} place du motif}\\[4pt]
100 &= 16\times 6+4 &&\justif{$16$ motifs complets ($96$ chiffres), puis $4$ rangs}\\
u_{100} &= u_{4} = \res{8} &&\justif{la 4\textsuperscript{e} place du motif}
\end{calculs}
\rem{Si le reste vaut $0$, on tombe sur la fin d'un motif : par exemple $u_{18}=u_{6}=7$.}
\end{exo}

\partie{Suites définies explicitement}

\begin{exo}[Calculer des termes]
\explique{Chaque suite est donnée par son terme général : pour obtenir un terme, on
remplace $n$ par le rang voulu.}{inutile de calculer les termes précédents}
\cas{a)}
\begin{calculs}
u_{0} &= 4\times 0-3 = \res{-3} &&\justif{on remplace $n$ par $0$}\\
u_{1} &= 4\times 1-3 = \res{1} &&\justif{on remplace $n$ par $1$}\\
u_{5} &= 4\times 5-3 &&\justif{on remplace $n$ par $5$}\\
u_{5} &= 20-3 = \res{17} &&
\end{calculs}
\cas{b)}
\begin{calculs}
v_{0} &= 0^{2}-2\times 0+3 = \res{3} &&\justif{on remplace $n$ par $0$}\\
v_{1} &= 1^{2}-2\times 1+3 &&\justif{on remplace $n$ par $1$}\\
v_{1} &= 1-2+3 = \res{2} &&\\
v_{5} &= 5^{2}-2\times 5+3 &&\justif{on remplace $n$ par $5$}\\
v_{5} &= 25-10+3 = \res{18} &&
\end{calculs}
\cas{c)}
\begin{calculs}
w_{0} &= \dfrac{3\times 0+2}{0+1} = \res{2} &&\justif{on remplace chaque $n$ par $0$}\\
w_{1} &= \dfrac{3\times 1+2}{1+1} = \res{\dfrac{5}{2}} &&\justif{on remplace chaque $n$
par $1$}\\
w_{5} &= \dfrac{3\times 5+2}{5+1} = \res{\dfrac{17}{6}} &&\justif{on remplace chaque
$n$ par $5$}
\end{calculs}
\rem{On garde la valeur exacte $\dfrac{17}{6}$ (irréductible : $17$ n'a pas d'autre
diviseur que $1$ et $17$) plutôt que l'arrondi $2{,}83$.}
\cas{d)}
\begin{calculs}
t_{0} &= 2^{0}-1 &&\justif{on remplace $n$ par $0$}\\
t_{0} &= 1-1 = \res{0} &&\justif{$2^{0}=1$, comme tout nombre non nul à la puissance
$0$}\\
t_{1} &= 2^{1}-1 = \res{1} &&\justif{$2^{1}=2$}\\
t_{5} &= 2^{5}-1 &&\justif{on remplace $n$ par $5$}\\
t_{5} &= 32-1 = \res{31} &&\justif{$2^{5}=2\times 2\times 2\times 2\times 2=32$}
\end{calculs}
\rem{Erreur fréquente : $2^{0}=0$ ou $2^{5}=10$. La puissance n'est pas un produit par
l'exposant.}
\end{exo}

\begin{exo}[Retrouver un rang]
\cas{1.}
\begin{calculs}
u_{10} &= 3\times 10+7 &&\justif{on remplace $n$ par $10$}\\
u_{10} &= 30+7 = \res{37} &&
\end{calculs}
\cas{2.}
\explique{On connaît le terme ($52$) et on cherche le rang : on résout l'équation
$u_{n}=52$, d'inconnue $n$.}{le chemin inverse de la question 1}
\begin{calculs}
3n+7 &= 52 &&\justif{on remplace $u_{n}$ par son expression}\\
3n &= 45 &&\justif{on retire $7$ aux deux membres}\\
n &= 15 &&\justif{on divise par $3$}
\end{calculs}
\explique{$15$ est un entier naturel : $\res{u_{15}=52}$, le terme $52$ est de rang
$15$.}{vérification : $3\times 15+7=45+7=52$}
\cas{3.}
\begin{calculs}
3n+7 &= 101 &&\justif{on résout $u_{n}=101$}\\
3n &= 94 &&\justif{on retire $7$}\\
n &= \dfrac{94}{3} \approx 31{,}3 &&\justif{on divise par $3$}
\end{calculs}
\explique{$\dfrac{94}{3}$ n'est pas un entier : aucun rang ne convient, donc
$\res{101 \text{ n'est pas un terme de la suite}}$.}{un rang est toujours un entier
naturel}
\rem{On le voit aussi en encadrant : $u_{31}=100$ et $u_{32}=103$ ; la suite saute
$101$.}
\end{exo}

\begin{exo}[Exprimer $u_{n+1}$ et $u_{n-1}$]
\explique{Pour obtenir $u_{n+1}$, on remplace $n$ par $(n+1)$ \emph{partout} dans
l'expression ; pour $u_{n-1}$, par $(n-1)$. Les parenthèses sont
indispensables.}{on change le rang, pas le terme}
\cas{a)}
\begin{calculs}
u_{n+1} &= 5(n+1)-2 &&\justif{on remplace $n$ par $(n+1)$}\\
u_{n+1} &= 5n+5-2 &&\justif{on développe}\\
u_{n+1} &= \res{5n+3} &&\\[4pt]
u_{n-1} &= 5(n-1)-2 &&\justif{on remplace $n$ par $(n-1)$}\\
u_{n-1} &= 5n-5-2 &&\\
u_{n-1} &= \res{5n-7} &&
\end{calculs}
\explique{Vérification avec $n=2$ : $u_{3}=5\times 3-2=13$ et $5\times 2+3=13$ ;
$u_{1}=5\times 1-2=3$ et $5\times 2-7=3$.}{on teste la formule sur une valeur}
\cas{b)}
\begin{calculs}
v_{n+1} &= (n+1)^{2}-4(n+1) &&\justif{on remplace $n$ par $(n+1)$}\\
v_{n+1} &= n^{2}+2n+1-4n-4 &&\justif{$(a+b)^{2}=a^{2}+2ab+b^{2}$}\\
v_{n+1} &= \res{n^{2}-2n-3} &&\justif{on regroupe : $2n-4n=-2n$ et $1-4=-3$}\\[4pt]
v_{n-1} &= (n-1)^{2}-4(n-1) &&\justif{on remplace $n$ par $(n-1)$}\\
v_{n-1} &= n^{2}-2n+1-4n+4 &&\justif{$(a-b)^{2}=a^{2}-2ab+b^{2}$ et $-4\times(-1)=+4$}\\
v_{n-1} &= \res{n^{2}-6n+5} &&
\end{calculs}
\explique{Vérification avec $n=3$ : $v_{4}=16-16=0$ et $9-6-3=0$ ;
$v_{2}=4-8=-4$ et $9-18+5=-4$.}{même test}
\cas{c)}
\begin{calculs}
w_{n+1} &= \dfrac{(n+1)+3}{2(n+1)+1} &&\justif{on remplace \emph{chaque} $n$ par
$(n+1)$}\\
w_{n+1} &= \res{\dfrac{n+4}{2n+3}} &&\justif{$2(n+1)+1=2n+2+1$}\\[4pt]
w_{n-1} &= \dfrac{(n-1)+3}{2(n-1)+1} &&\justif{on remplace chaque $n$ par $(n-1)$}\\
w_{n-1} &= \res{\dfrac{n+2}{2n-1}} &&\justif{$2(n-1)+1=2n-2+1$}
\end{calculs}
\rem{Erreur fréquente : ne remplacer $n$ qu'au numérateur. Pour $n\ge 1$, le dénominateur
$2n-1$ vaut au moins $1$ : $w_{n-1}$ existe bien.}
\end{exo}

\begin{exo}[Ne pas confondre]
\explique{Deux gestes différents : changer le rang (on remplace $n$ dans l'expression),
ou opérer sur le terme (on garde $u_{n}$ et on calcule avec).}{le cœur de l'exercice}
\cas{a)}
\begin{calculs}
u_{n+1} &= (n+1)^{2}+2(n+1) &&\justif{on remplace $n$ par $(n+1)$ : terme suivant}\\
u_{n+1} &= n^{2}+2n+1+2n+2 &&\justif{$(n+1)^{2}=n^{2}+2n+1$}\\
u_{n+1} &= \res{n^{2}+4n+3} &&
\end{calculs}
\cas{b)}
\begin{calculs}
u_{n}+1 &= n^{2}+2n+1 &&\justif{on ajoute $1$ au terme $u_{n}$ ; $n$ ne change pas}
\end{calculs}
\explique{$\res{u_{n}+1=n^{2}+2n+1}$.}{à comparer avec a)}
\cas{c)}
\begin{calculs}
u_{2n} &= (2n)^{2}+2\times(2n) &&\justif{on remplace $n$ par $(2n)$ : terme de rang
double}\\
u_{2n} &= \res{4n^{2}+4n} &&\justif{$(2n)^{2}=2^{2}\times n^{2}=4n^{2}$}
\end{calculs}
\cas{d)}
\begin{calculs}
2u_{n} &= 2(n^{2}+2n) &&\justif{on multiplie le terme $u_{n}$ par $2$}\\
2u_{n} &= \res{2n^{2}+4n} &&
\end{calculs}
\explique{Quatre expressions différentes. Vérification avec $n=1$ : $u_{1}=3$ et
$u_{2}=8$ ; on trouve bien $u_{n+1}=1+4+3=8$, alors que $u_{n}+1=4$ ; de même
$u_{2n}=u_{2}=8$, alors que $2u_{n}=6$.}{un test numérique sépare les quatre}
\rem{À retenir : un calcul \emph{dans l'indice} ($n+1$ ou $2n$) change le rang ; un calcul
\emph{sur} $u_{n}$ ($+1$, $\times 2$) change la valeur du terme.}
\end{exo}

\partie{Suites définies par récurrence}

\begin{exo}[Calculer de proche en proche]
\explique{On applique la relation de récurrence une fois par terme : avec $n=0$, elle
donne $u_{1}$ à partir de $u_{0}$ ; avec $n=1$, $u_{2}$ à partir de $u_{1}$\dots}{chaque
terme a besoin du précédent}
\cas{a)}
\begin{calculs}
u_{1} &= 2u_{0}-1 &&\justif{la relation avec $n=0$}\\
u_{1} &= 2\times 3-1 = 5 &&\justif{on remplace $u_{0}$ par sa valeur $3$}\\[4pt]
u_{2} &= 2u_{1}-1 &&\justif{la relation avec $n=1$}\\
u_{2} &= 2\times 5-1 = 9 &&\justif{$u_{1}=5$}\\[4pt]
u_{3} &= 2u_{2}-1 &&\justif{la relation avec $n=2$}\\
u_{3} &= 2\times 9-1 = \res{17} &&
\end{calculs}
\explique{Les trois termes qui suivent $u_{0}$ : $\res{u_{1}=5}$ ; $\res{u_{2}=9}$ ;
$\res{u_{3}=17}$.}{}
\cas{b)}
\explique{Ici la relation contient aussi $n$ : on remplace $n$ par le rang du terme
\emph{dont on part}.}{$n=0$ pour calculer $v_{1}$}
\begin{calculs}
v_{1} &= v_{0}+3\times 0-1 &&\justif{$n=0$ : on remplace aussi $n$ par $0$}\\
v_{1} &= 2+0-1 = 1 &&\\[4pt]
v_{2} &= v_{1}+3\times 1-1 &&\justif{$n=1$ : le rang du terme précédent}\\
v_{2} &= 1+3-1 = 3 &&\\[4pt]
v_{3} &= v_{2}+3\times 2-1 &&\justif{$n=2$}\\
v_{3} &= 3+6-1 = \res{8} &&
\end{calculs}
\explique{$\res{v_{1}=1}$ ; $\res{v_{2}=3}$ ; $\res{v_{3}=8}$.}{}
\rem{Erreur fréquente : écrire $v_{2}=v_{1}+3\times 2-1$. Pour obtenir $v_{n+1}$, on
prend $n$, le rang du terme précédent : pour $v_{2}$, c'est $n=1$.}
\cas{c)}
\explique{Le premier terme est $w_{1}$ : les trois suivants sont $w_{2}$, $w_{3}$ et
$w_{4}$.}{on part de $n=1$}
\begin{calculs}
w_{2} &= \dfrac{w_{1}}{2}+4 &&\justif{la relation avec $n=1$}\\
w_{2} &= \dfrac{16}{2}+4 = 12 &&\justif{$w_{1}=16$}\\[4pt]
w_{3} &= \dfrac{12}{2}+4 = 10 &&\justif{$n=2$}\\
w_{4} &= \dfrac{10}{2}+4 = \res{9} &&\justif{$n=3$}
\end{calculs}
\explique{$\res{w_{2}=12}$ ; $\res{w_{3}=10}$ ; $\res{w_{4}=9}$.}{}
\end{exo}

\begin{exo}[Avec un carré, avec une fraction]
\cas{a)}
\explique{$a_{n}^{2}$ est le carré du \emph{terme} $a_{n}$, pas celui du rang.}{on élève
au carré le terme précédent}
\begin{calculs}
a_{1} &= a_{0}^{2}+1 &&\justif{la relation avec $n=0$}\\
a_{1} &= 1^{2}+1 = 2 &&\justif{$a_{0}=1$}\\[4pt]
a_{2} &= 2^{2}+1 = 5 &&\justif{$a_{1}=2$}\\
a_{3} &= 5^{2}+1 = \res{26} &&\justif{$a_{2}=5$}
\end{calculs}
\explique{$\res{a_{1}=2}$ ; $\res{a_{2}=5}$ ; $\res{a_{3}=26}$.}{}
\cas{b)}
\begin{calculs}
b_{1} &= \dfrac{1}{1+b_{0}} &&\justif{la relation avec $n=0$}\\
b_{1} &= \dfrac{1}{1+1} = \dfrac{1}{2} &&\justif{$b_{0}=1$}\\[4pt]
b_{2} &= \dfrac{1}{1+\frac{1}{2}} &&\justif{$b_{1}=\frac{1}{2}$}\\
b_{2} &= \dfrac{1}{\frac{2}{2}+\frac{1}{2}} &&\justif{$1=\frac{2}{2}$ : même
dénominateur}\\
b_{2} &= \dfrac{1}{\frac{3}{2}} = \dfrac{2}{3} &&\justif{diviser par $\frac{3}{2}$,
c'est multiplier par $\frac{2}{3}$}\\[4pt]
b_{3} &= \dfrac{1}{1+\frac{2}{3}} &&\justif{$b_{2}=\frac{2}{3}$}\\
b_{3} &= \dfrac{1}{\frac{3}{3}+\frac{2}{3}} &&\justif{$1=\frac{3}{3}$}\\
b_{3} &= \dfrac{1}{\frac{5}{3}} = \res{\dfrac{3}{5}} &&\justif{l'inverse de
$\frac{5}{3}$}
\end{calculs}
\explique{$\res{b_{1}=\frac{1}{2}}$ ; $\res{b_{2}=\frac{2}{3}}$ ;
$\res{b_{3}=\frac{3}{5}}$.}{vérification à la calculatrice : $0{,}5$ ; $0{,}666\dots$ ;
$0{,}6$}
\end{exo}

\begin{exo}[Traduire une situation]
\explique{Pour chaque situation : $u_{0}$ est la valeur de départ ($n=0$) ; la relation
dit comment on passe d'une étape ($u_{n}$) à la suivante ($u_{n+1}$).}{premier terme et
relation de récurrence}
\cas{a)}
\explique{$\res{u_{0}=1\,800}$.}{le salaire de départ}
\explique{Augmenter de $2\,\%$, c'est multiplier par $1+0{,}02=1{,}02$ :
$\res{u_{n+1}=1{,}02\,u_{n}}$.}{coefficient multiplicateur}
\explique{Vérification : $u_{1}=1{,}02\times 1\,800=1\,836$, soit $1\,800+36$.}{$2\,\%$
de $1\,800$ : $36$}
\cas{b)}
\explique{$\res{u_{0}=12\,000}$.}{la population de départ}
\explique{Si $5\,\%$ des habitants partent, il en reste $95\,\%$ : on multiplie par
$1-0{,}05=0{,}95$. Puis on ajoute les $900$ arrivants :
$\res{u_{n+1}=0{,}95\,u_{n}+900}$.}{d'abord le pourcentage, puis l'ajout}
\rem{Le pourcentage porte sur la population de l'année $n$ : on multiplie $u_{n}$, puis
on ajoute $900$. Écrire $0{,}95(u_{n}+900)$ ferait partir $5\,\%$ des nouveaux
arrivants.}
\cas{c)}
\explique{$\res{u_{0}=500}$ et $\res{u_{n+1}=u_{n}-40}$.}{on retire chaque jour la même
quantité}
\rem{Perdre $40$~L, c'est soustraire $40$ ; perdre $40\,\%$, ce serait multiplier par
$0{,}6$. Bien lire s'il s'agit d'une quantité ou d'un pourcentage.}
\end{exo}

\begin{exo}[Une forêt]
\cas{1.}
\explique{$u_{0}=5\,000$.}{en 2026 : $n=0$}
\begin{calculs}
u_{1} &= 5\,000-0{,}08\times 5\,000+300 &&\justif{on abat $8\,\%$, on plante $300$}\\
u_{1} &= 5\,000-400+300 = \res{4\,900} &&\justif{$8\,\%$ de $5\,000$ : $400$}\\[4pt]
u_{2} &= 4\,900-0{,}08\times 4\,900+300 &&\justif{même calcul à partir de $u_{1}$}\\
u_{2} &= 4\,900-392+300 = \res{4\,808} &&\justif{$8\,\%$ de $4\,900$ : $392$}
\end{calculs}
\explique{La forêt compte $4\,900$ arbres en 2027 et $4\,808$ en 2028.}{on répond avec
les années}
\cas{2.}
\explique{Abattre $8\,\%$ des arbres, c'est multiplier par $1-0{,}08=0{,}92$ : il reste
$0{,}92\,u_{n}$ arbres.}{coefficient multiplicateur}
\explique{On ajoute ensuite les $300$ arbres plantés :
$\res{u_{n+1}=0{,}92\,u_{n}+300}$.}{d'une année à la suivante}
\explique{Vérification : $0{,}92\times 4\,900+300=4\,508+300=4\,808$, c'est bien
$u_{2}$.}{la relation redonne la question 1}
\end{exo}

\partie{Calculatrice}

\begin{exo}[Afficher des termes]
\cas{1.}
\explique{\texttt{Explicite} : $u_{n}=n^{2}-10n+30$ ; \texttt{Tableau} : \texttt{N
début} $0$, \texttt{N fin} $10$.}{du rang $0$ au rang $10$}
\begin{center}
\renewcommand{\arraystretch}{1.3}
\begin{tabular}{|c|*{11}{c|}}
\hline
$n$ & $0$ & $1$ & $2$ & $3$ & $4$ & $5$ & $6$ & $7$ & $8$ & $9$ & $10$\\\hline
$u_{n}$ & $30$ & $21$ & $14$ & $9$ & $6$ & $5$ & $6$ & $9$ & $14$ & $21$ & $30$\\\hline
\end{tabular}
\end{center}
\explique{Vérification d'une valeur à la main : $u_{3}=9-30+30=9$.}{pour s'assurer que
la suite est bien entrée}
\cas{2.}
\explique{\texttt{Graphique}, \texttt{Axes} : par exemple $X_{\min}=0$, $X_{\max}=11$,
$Y_{\min}=0$, $Y_{\max}=32$.}{un peu au-delà des valeurs du tableau : $n$ va jusqu'à
$10$, $u_{n}$ reste entre $5$ et $30$}
\explique{Les points descendent jusqu'au rang $5$, puis remontent.}{ce qu'on voit à
l'écran}
\cas{3.}
\explique{\texttt{Récurrente d'ordre 1} : $v_{n+1}=2v_{n}-2$ et $v_{0}=3$.}{la relation
et le premier terme}
\explique{\texttt{Tableau} jusqu'à $20$ : $\res{v_{10}=1\,026}$ et
$\res{v_{20}=1\,048\,578}$.}{la calculatrice enchaîne les termes à notre place}
\rem{Pour contrôler la saisie, on compare les premiers termes à un calcul à la main :
$v_{1}=2\times 3-2=4$, puis $v_{2}=6$ et $v_{3}=10$. Une erreur fréquente est un premier
terme mal entré ($v_{1}$ au lieu de $v_{0}$).}
\end{exo}

\begin{exo}[Conjecturer une expression]
\cas{1.}
\explique{\texttt{Récurrente d'ordre 1} : $w_{n+1}=w_{n}+2n+1$ et $w_{0}=1$, puis
\texttt{Tableau}.}{$n$ s'obtient avec la touche \fbox{\small$x,n,t$}}
\begin{center}
\renewcommand{\arraystretch}{1.3}
\begin{tabular}{|c|*{7}{c|}}
\hline
$n$ & $0$ & $1$ & $2$ & $3$ & $4$ & $5$ & $6$\\\hline
$w_{n}$ & $1$ & $2$ & $5$ & $10$ & $17$ & $26$ & $37$\\\hline
\end{tabular}
\end{center}
\explique{Plus loin dans le tableau : $\res{w_{10}=101}$ et
$\res{w_{50}=2\,501}$.}{contrôle à la main : $w_{1}=1+2\times 0+1=2$}
\cas{2.}
\explique{On retire $1$ à chaque terme : $0$ ; $1$ ; $4$ ; $9$ ; $16$ ; $25$ ; $36$, ce
sont les carrés $0^{2}$ ; $1^{2}$ ; $2^{2}$\dots{} des rangs.}{on cherche un lien entre
le terme et son rang}
\explique{On conjecture que, pour tout $n$, $\res{w_{n}=n^{2}+1}$.}{une conjecture,
qu'on ne démontre pas ici}
\explique{Elle s'accorde avec les termes lointains : $10^{2}+1=101$ et
$50^{2}+1=2\,501$.}{on vérifie sur les valeurs de la question 1}
\end{exo}

\begin{exo}[Représenter et conjecturer]
\cas{1.}
\explique{\texttt{Explicite} : $t_{n}=\dfrac{20}{n+1}$ ; \texttt{Tableau} de $0$ à
$9$.}{arrondis à $0{,}01$}
\begin{center}
\renewcommand{\arraystretch}{1.3}
\begin{tabular}{|c|*{10}{c|}}
\hline
$n$ & $0$ & $1$ & $2$ & $3$ & $4$ & $5$ & $6$ & $7$ & $8$ & $9$\\\hline
$t_{n}$ & $20$ & $10$ & $6{,}67$ & $5$ & $4$ & $3{,}33$ & $2{,}86$ & $2{,}5$ & $2{,}22$ &
$2$\\\hline
\end{tabular}
\end{center}
\cas{2.}
\explique{\texttt{Graphique}, \texttt{Axes} : par exemple $X$ de $0$ à $10$ et $Y$ de
$0$ à $22$.}{fenêtre qui contient tous les points}
\explique{Les points descendent de gauche à droite : la suite semble
\res{décroissante}.}{une conjecture, lue sur le graphique}
\explique{On le comprend : le dénominateur $n+1$ augmente, donc la fraction
$\dfrac{20}{n+1}$ diminue.}{$20$ partagé en de plus en plus de parts}
\end{exo}

\begin{exo}[Le carré ou la puissance ?]
\cas{1.}
\explique{Deux suites \texttt{Explicite} : $u_{n}=n^{2}$ et $v_{n}=2^{n}$ ; le
\texttt{Tableau} les affiche côte à côte.}{une colonne par suite}
\begin{center}
\renewcommand{\arraystretch}{1.3}
\begin{tabular}{|c|*{11}{c|}}
\hline
$n$ & $0$ & $1$ & $2$ & $3$ & $4$ & $5$ & $6$ & $7$ & $8$ & $9$ & $10$\\\hline
$u_{n}=n^{2}$ & $0$ & $1$ & $4$ & $9$ & $16$ & $25$ & $36$ & $49$ & $64$ & $81$ &
$100$\\\hline
$v_{n}=2^{n}$ & $1$ & $2$ & $4$ & $8$ & $16$ & $32$ & $64$ & $128$ & $256$ & $512$ &
$1\,024$\\\hline
\end{tabular}
\end{center}
\cas{2.}
\explique{On compare les deux lignes colonne par colonne : $u_{n}<v_{n}$ pour $n=0$ et
$n=1$ ; $u_{2}=v_{2}=4$ ; $u_{3}=9>8=v_{3}$ ; $u_{4}=v_{4}=16$ ; puis $u_{n}<v_{n}$
jusqu'à $n=10$.}{une comparaison par rang}
\explique{$u_{n}>v_{n}$ seulement pour $\res{n=3}$ ; $u_{n}=v_{n}$ pour $\res{n=2}$ et
$\res{n=4}$.}{d'après le tableau}
\cas{3.}
\explique{On conjecture que, pour tout $n\ge 5$, $\res{2^{n}>n^{2}}$.}{le tableau le
montre jusqu'à $n=10$}
\explique{L'écart grandit vite : $32-25=7$, puis $64-36=28$, \dots,
$1\,024-100=924$.}{il ne semble pas pouvoir se refermer}
\rem{D'un rang au suivant, $2^{n}$ double, alors que $n^{2}$ augmente beaucoup moins :
de $100$ à $121$ entre les rangs $10$ et $11$.}
\end{exo}

\partie{Programmer en Python}

\begin{exo}[Écrire la fonction]
\cas{1.}
\begin{minipage}[t]{0.42\linewidth}\vspace{0pt}
\begin{python}
def terme_v(n):
    v = 4
    for i in range(n):
        v = 3*v - 5
    return v
\end{python}
\end{minipage}\hfill
\begin{minipage}[t]{0.54\linewidth}\vspace{0pt}\raggedright
\gris{\texttt{v = 4} donne $v_{0}$ ; la boucle tourne $n$ fois, de $v_{0}$ à $v_{n}$ ;
dans \texttt{v = 3*v - 5}, le \texttt{v} de droite est le terme précédent, celui de
gauche le terme suivant : c'est la relation $v_{n+1}=3v_{n}-5$}
\end{minipage}
\rem{Erreurs fréquentes : placer \texttt{v = 4} dans la boucle (on repartirait de $v_{0}$
à chaque tour) ; décaler \texttt{return v} dans la boucle (la fonction s'arrêterait au
premier tour).}
\cas{2.}
\explique{\texttt{terme\_v(10)} renvoie \texttt{88576} : $\res{v_{10}=88\,576}$.}{dans la
console : \texttt{terme\_v(10)}}
\explique{Contrôle : \texttt{terme\_v(1)} renvoie $7$, et à la main
$v_{1}=3\times 4-5=7$.}{un petit rang se vérifie à la main}
\end{exo}

\begin{exo}[Lire un programme]
\cas{1.}
\explique{\texttt{u = 5} : le premier terme est $\res{u_{0}=5}$.}{la valeur avant la
boucle}
\explique{\texttt{u = u/2 + 3} : à chaque tour, le nouveau terme est la moitié du
précédent, plus $3$ : $\res{u_{n+1}=\dfrac{u_{n}}{2}+3}$.}{la ligne de la boucle est la
relation de récurrence}
\explique{La boucle tourne $n$ fois : \texttt{mystere(n)} renvoie $u_{n}$.}{de $u_{0}$
à $u_{n}$}
\cas{2.}
\begin{calculs}
u_{1} &= \dfrac{u_{0}}{2}+3 &&\justif{un tour de boucle}\\
u_{1} &= \dfrac{5}{2}+3 = 5{,}5 &&\justif{$\frac{5}{2}=2{,}5$}\\[4pt]
u_{2} &= \dfrac{5{,}5}{2}+3 = 5{,}75 &&\justif{deux tours}\\
u_{3} &= \dfrac{5{,}75}{2}+3 = 5{,}875 &&\justif{trois tours}
\end{calculs}
\explique{\texttt{mystere(1)} renvoie $\res{5{,}5}$ ; \texttt{mystere(2)} renvoie
$\res{5{,}75}$ ; \texttt{mystere(3)} renvoie $\res{5{,}875}$.}{on conclut sur ce que
renvoie la fonction}
\end{exo}

\begin{exo}[Une relation qui dépend de $n$]
\cas{1.}
\begin{calculs}
w_{1} &= w_{0}+2\times 0+3 &&\justif{la relation avec $n=0$}\\
w_{1} &= 0+0+3 = \res{3} &&\\[4pt]
w_{2} &= w_{1}+2\times 1+3 &&\justif{$n=1$}\\
w_{2} &= 3+2+3 = \res{8} &&\\[4pt]
w_{3} &= w_{2}+2\times 2+3 &&\justif{$n=2$}\\
w_{3} &= 8+4+3 = \res{15} &&
\end{calculs}
\cas{2.}
\begin{minipage}[t]{0.42\linewidth}\vspace{0pt}
\begin{python}
def terme_w(n):
    w = 0
    for i in range(n):
        w = w + 2*i + 3
    return w
\end{python}
\end{minipage}\hfill
\begin{minipage}[t]{0.54\linewidth}\vspace{0pt}\raggedright
\gris{\texttt{i} prend les valeurs $0$ ; $1$ ; \dots ; $n-1$ ; au tour où \texttt{i} vaut
$k$, \texttt{w} contient $w_{k}$ et devient $w_{k+1}=w_{k}+2k+3$ : \texttt{i} joue le
rôle du $n$ de la relation}
\end{minipage}
\rem{Erreur fréquente : écrire \texttt{2*n} dans la boucle. Le paramètre \texttt{n} est
le nombre de tours, il ne change pas ; c'est \texttt{i} qui avance d'un rang à chaque
tour.}
\cas{3.}
\explique{\texttt{terme\_w(10)} renvoie $\res{120}$ : $w_{10}=120$.}{contrôle :
\texttt{terme\_w(3)} renvoie $15$, comme à la question 1}
\end{exo}

\begin{exo}[La suite de Fibonacci]
\cas{1.}
\explique{Chaque terme est la somme des deux précédents : $f_{n+2}=f_{n+1}+f_{n}$.}{il
faut deux termes pour avancer}
\begin{calculs}
f_{2} &= f_{1}+f_{0} &&\justif{la relation avec $n=0$}\\
f_{2} &= 1+1 = 2 &&\\
f_{3} &= 2+1 = 3 &&\justif{$f_{2}+f_{1}$}\\
f_{4} &= 3+2 = 5 &&\justif{$f_{3}+f_{2}$}\\
f_{5} &= 5+3 = 8 &&\justif{$f_{4}+f_{3}$}\\
f_{6} &= 8+5 = \res{13} &&\justif{$f_{5}+f_{4}$}
\end{calculs}
\cas{2.}
\explique{Ligne à compléter : \texttt{a, b = b, a + b}.}{le nouveau terme est la somme
des deux précédents}
\explique{Si \texttt{a} et \texttt{b} contiennent $f_{k}$ et $f_{k+1}$, ils reçoivent
$f_{k+1}$ et $f_{k}+f_{k+1}=f_{k+2}$ : on avance d'un rang.}{les deux valeurs de droite
sont calculées avec les anciennes valeurs}
\explique{Au départ, \texttt{b} contient $f_{1}$ ; après $n-1$ tours, il contient
$f_{n}$.}{d'où \texttt{range(n-1)}}
\begin{center}
\renewcommand{\arraystretch}{1.3}
\begin{tabular}{|c|c|c|c|c|}
\hline
tour & départ & $1$ & $2$ & $3$\\\hline
\texttt{a} & $1$ & $1$ & $2$ & $3$\\\hline
\texttt{b} & $1$ & $2$ & $3$ & $5$\\\hline
\end{tabular}
\end{center}
\explique{Trois tours pour \texttt{fibo(4)} : la fonction renvoie $5=f_{4}$.}{on suit
les variables pour contrôler}
\rem{Écrire \texttt{a = b} puis \texttt{b = a + b} sur deux lignes est faux : \texttt{a}
a déjà changé, et \texttt{b} serait seulement doublé.}
\cas{3.}
\explique{\texttt{fibo(20)} renvoie $\res{10\,946}$ : $f_{20}=10\,946$.}{dans la
console}
\end{exo}

\partie{Représentation graphique}

\begin{exo}[Représenter une suite récurrente]
\cas{1.}
\begin{calculs}
u_{1} &= 0{,}5\times 2+2 = 3 &&\justif{la relation avec $n=0$}\\
u_{2} &= 0{,}5\times 3+2 = 3{,}5 &&\justif{$n=1$}\\
u_{3} &= 0{,}5\times 3{,}5+2 = 3{,}75 &&\justif{$n=2$}\\
u_{4} &= 0{,}5\times 3{,}75+2 = 3{,}875 &&\justif{$n=3$}\\
u_{5} &= 0{,}5\times 3{,}875+2 = 3{,}9375 &&\justif{$n=4$}
\end{calculs}
\begin{center}
\renewcommand{\arraystretch}{1.3}
\begin{tabular}{|c|*{6}{c|}}
\hline
$n$ & $0$ & $1$ & $2$ & $3$ & $4$ & $5$\\\hline
$u_{n}$ & $2$ & $3$ & $3{,}5$ & $3{,}75$ & $3{,}875$ & $3{,}9375$\\\hline
\end{tabular}
\end{center}
\cas{2. et 3.}
\begin{minipage}[c]{0.58\linewidth}\raggedright
On place un point par colonne du tableau : $(0\,;2)$, $(1\,;3)$, $(2\,;3{,}5)$\dots{}
Les points restent isolés : on ne les relie pas.

\smallskip
Les points montent de gauche à droite : la suite semble
\res{croissante}, et semble se rapprocher de $4$.

\smallskip
\gris{L'écart à $4$ est divisé par $2$ à chaque rang : $2$, puis $1$, puis $0{,}5$\dots{}
C'est une conjecture : le sens de variation se démontre au chapitre 7.}
\end{minipage}\hfill
\begin{minipage}[c]{0.38\linewidth}\centering
\begin{tikzpicture}[x=0.6cm,y=0.6cm]
  \reperesuite{6}{0}{5}{u_{n}}
  \gradx{1,2,3,4,5} \grady{1,2,3,4}
  \foreach \n/\u in {0/2,1/3,2/3.5,3/3.75,4/3.875,5/3.9375}
    \path[coulCourbe] (\n,\u) node[croix]{};
\end{tikzpicture}
\end{minipage}
\end{exo}

\begin{exo}[Représenter une suite explicite]
\cas{1.}
\begin{calculs}
v_{0} &= \dfrac{(0-3)^{2}}{2} = \dfrac{9}{2} &&\justif{on remplace $n$ par $0$ ;
$(-3)^{2}=9$}\\
v_{0} &= 4{,}5 &&\\
v_{1} &= \dfrac{(1-3)^{2}}{2} = 2 &&\justif{$(-2)^{2}=4$}\\
v_{2} &= \dfrac{(2-3)^{2}}{2} = 0{,}5 &&\justif{$(-1)^{2}=1$}\\
v_{3} &= \dfrac{(3-3)^{2}}{2} = 0 &&
\end{calculs}
\explique{Pour $n=4$ ; $5$ ; $6$, on retrouve les mêmes carrés $1$ ; $4$ ; $9$ :
$v_{4}=0{,}5$ ; $v_{5}=2$ ; $v_{6}=4{,}5$.}{$(n-3)^{2}$ ne dépend que de l'écart entre
$n$ et $3$}
\begin{center}
\renewcommand{\arraystretch}{1.3}
\begin{tabular}{|c|*{7}{c|}}
\hline
$n$ & $0$ & $1$ & $2$ & $3$ & $4$ & $5$ & $6$\\\hline
$v_{n}$ & $4{,}5$ & $2$ & $0{,}5$ & $0$ & $0{,}5$ & $2$ & $4{,}5$\\\hline
\end{tabular}
\end{center}
\cas{2. et 3.}
\begin{minipage}[c]{0.58\linewidth}\raggedright
On place les points $(n\,;v_{n})$ du tableau.

\smallskip
Les points descendent de $v_{0}$ à $v_{3}$, puis remontent : la suite semble
\res{décroissante jusqu'au rang $3$}, puis
\res{croissante à partir du rang $3$}.

\smallskip
\gris{Le plus petit terme est $v_{3}=0$ ; la figure est symétrique par rapport au rang
$3$.}
\end{minipage}\hfill
\begin{minipage}[c]{0.38\linewidth}\centering
\begin{tikzpicture}[x=0.5cm,y=0.5cm]
  \reperesuite{7}{0}{5}{v_{n}}
  \gradx{1,2,3,4,5,6} \grady{1,2,3,4}
  \foreach \n in {0,1,...,6}
    \path[coulCourbe] (\n,{(\n-3)^2/2}) node[croix]{};
\end{tikzpicture}
\end{minipage}
\end{exo}

\begin{exo}[Lire un nuage de points]
\begin{minipage}[c]{0.58\linewidth}\raggedright
\cas{1.}
\explique{Le point d'abscisse $n$ a pour ordonnée $u_{n}$.}{abscisse : le rang ;
ordonnée : le terme}
\explique{$\res{u_{0}=1}$ ; $\res{u_{2}=4}$ ; $\res{u_{5}=1}$.}{on lit l'ordonnée au-dessus
de $0$ ; $2$ et $5$}
\end{minipage}\hfill
\begin{minipage}[c]{0.38\linewidth}\centering
\begin{tikzpicture}[x=0.45cm,y=0.45cm]
  \reperesuite{7}{0}{5}{u_{n}}
  \gradx{1,2,3,4,5,6} \grady{1,2,3,4}
  \foreach \n/\u in {0/1,1/3,2/4,3/4,4/3,5/1,6/0}
    \path[coulCourbe] (\n,\u) node[croix]{};
\end{tikzpicture}
\end{minipage}
\cas{2.}
\explique{Les points d'ordonnée $4$ ont pour abscisses $2$ et $3$ : $u_{n}=4$ pour
$\res{n=2}$ et $\res{n=3}$.}{on lit cette fois des abscisses}
\rem{La question demande des rangs : répondre « $4$ » ou « $u_{4}$ » est une erreur
fréquente.}
\cas{3.}
\explique{$u_{0}<u_{1}<u_{2}$ : la suite \res{croît de $u_{0}$ à $u_{2}$}.}{$1$, puis $3$, puis $4$}
\explique{$u_{2}=u_{3}=4$, puis $u_{3}>u_{4}>u_{5}>u_{6}$ : elle
\res{décroît de $u_{3}$ à $u_{6}$}.}{$4$, puis $3$ ; $1$ et $0$}
\explique{Sur ces premiers termes, la suite n'est donc ni croissante ni
décroissante.}{on ne dit rien des termes au-delà de $u_{6}$}
\end{exo}

\begin{exo}[Des points sur deux droites]
\cas{1.}
\explique{$(-1)^{n}$ vaut $1$ si $n$ est pair, $-1$ si $n$ est impair : $u_{n}=n$ pour
$n$ pair, $u_{n}=-n$ pour $n$ impair.}{le signe alterne}
\begin{center}
\renewcommand{\arraystretch}{1.3}
\begin{tabular}{|c|*{7}{c|}}
\hline
$n$ & $0$ & $1$ & $2$ & $3$ & $4$ & $5$ & $6$\\\hline
$u_{n}$ & $0$ & $-1$ & $2$ & $-3$ & $4$ & $-5$ & $6$\\\hline
\end{tabular}
\end{center}
\begin{minipage}[c]{0.58\linewidth}\raggedright
\cas{2.}
\explique{$u_{1}<u_{0}$ ($-1<0$) : la suite n'est pas croissante.}{un seul contre-exemple
suffit}
\explique{$u_{2}>u_{1}$ ($2>-1$) : elle n'est pas décroissante.}{de même}
\explique{La suite n'est \res{ni croissante ni décroissante}.}{elle monte et
descend à chaque rang}
\end{minipage}\hfill
\begin{minipage}[c]{0.38\linewidth}\centering
\begin{tikzpicture}[x=0.32cm,y=0.32cm]
  \reperesuite{7}{-6}{7}{u_{n}}
  \gradx{2,4,6} \grady{-4,-2,2,4,6}
  \draw[coulCourbe!50,dashed] (0,0)--(6.8,6.8) (0,0)--(6.8,-6.8);
  \foreach \n in {0,1,...,6}
    \path[coulCourbe] (\n,{(-1)^\n*\n}) node[croix]{};
\end{tikzpicture}
\end{minipage}
\cas{3.}
\explique{Rangs pairs : $u_{n}=n$, les points $(n\,;n)$ sont sur la droite d'équation
$\res{y=x}$.}{l'ordonnée est égale à l'abscisse}
\explique{Rangs impairs : $u_{n}=-n$, les points $(n\,;-n)$ sont sur la droite
d'équation $\res{y=-x}$.}{l'ordonnée est l'opposée de l'abscisse}
\end{exo}

\partie{Sens de variation : signe de $u_{n+1}-u_{n}$}

\begin{exo}[Étudier le signe de la différence]
\explique{Comparer $u_{n+1}$ et $u_{n}$, c'est étudier le signe de leur différence : si
$u_{n+1}-u_{n}\ge 0$ pour tout $n$, alors $u_{n+1}\ge u_{n}$ et la suite est
croissante.}{propriété 1 : on remplace une comparaison par un signe}
\cas{a)}
\explique{On exprime d'abord $u_{n+1}$ :}{on remplace $n$ par $(n+1)$}
\begin{calculs}
u_{n+1} &= 3(n+1)^{2}+(n+1) &&\justif{on recopie l'expression}\\
u_{n+1} &= 3(n^{2}+2n+1)+n+1 &&\justif{$(n+1)^{2}=n^{2}+2n+1$}\\
u_{n+1} &= 3n^{2}+6n+3+n+1 &&\justif{on distribue le $3$}\\
u_{n+1} &= 3n^{2}+7n+4 &&
\end{calculs}
\explique{Puis la différence :}{on soustrait $u_{n}$ entre parenthèses}
\begin{calculs}
u_{n+1}-u_{n} &= 3n^{2}+7n+4-(3n^{2}+n) &&\justif{on soustrait $u_{n}$, entre parenthèses}\\
u_{n+1}-u_{n} &= 3n^{2}+7n+4-3n^{2}-n &&\justif{le signe $-$ change tous les signes}\\
u_{n+1}-u_{n} &= 6n+4 &&\justif{les $n^{2}$ s'éliminent}
\end{calculs}
\explique{Comme $n\ge 0$, on a $6n+4\ge 4>0$ : pour tout $n$, $u_{n+1}>u_{n}$. La suite
$(u_{n})$ est \res{strictement croissante}.}{une différence positive pour tout
$n$}
\explique{Vérification : $u_{0}=0$ et $u_{1}=4$, donc $u_{1}-u_{0}=4=6\times 0+4$.}{la
formule sur le premier rang}
\rem{Erreur fréquente : oublier les parenthèses autour de $u_{n}$ et écrire
$3n^{2}+7n+4-3n^{2}+n$.}
\cas{b)}
\begin{calculs}
v_{n+1}-v_{n} &= v_{n}+n^{2}+2-v_{n} &&\justif{on remplace $v_{n+1}$ par la relation}\\
v_{n+1}-v_{n} &= n^{2}+2 &&\justif{les $v_{n}$ s'éliminent}
\end{calculs}
\explique{$n^{2}\ge 0$, donc $n^{2}+2\ge 2>0$ : la suite $(v_{n})$ est
\res{strictement croissante}.}{pas besoin de connaître $v_{n}$ ni $v_{0}$}
\rem{Pour une suite définie par récurrence, la relation donne directement $v_{n+1}$ :
il n'y a rien à développer.}
\cas{c)}
\begin{calculs}
w_{n+1} &= (n+1)^{2}-10(n+1) &&\justif{on recopie l'expression}\\
w_{n+1} &= n^{2}+2n+1-10n-10 &&\justif{$(n+1)^{2}=n^{2}+2n+1$}\\
w_{n+1} &= n^{2}-8n-9 &&\\[4pt]
w_{n+1}-w_{n} &= n^{2}-8n-9-(n^{2}-10n) &&\justif{on soustrait $w_{n}$, entre parenthèses}\\
w_{n+1}-w_{n} &= n^{2}-8n-9-n^{2}+10n &&\justif{$-(-10n)=+10n$}\\
w_{n+1}-w_{n} &= 2n-9 &&
\end{calculs}
\explique{Le signe de $2n-9$ dépend de $n$ : on résout $2n-9\ge 0$.}{une inéquation du
premier degré}
\begin{calculs}
2n-9 &\ge 0 &&\\
2n &\ge 9 &&\justif{on ajoute $9$}\\
n &\ge 4{,}5 &&\justif{on divise par $2>0$}
\end{calculs}
\explique{$n$ est entier : pour $n\ge 5$, $w_{n+1}-w_{n}>0$ ; pour $n\le 4$, la
différence est négative.}{après $4{,}5$, le premier entier est $5$}
\explique{La suite $(w_{n})$ est \res{strictement croissante à partir du rang $5$} ; avant, elle décroît.}{$w_{4}=-24$ ; $w_{5}=-25$ ; $w_{6}=-24$}
\end{exo}

\begin{exo}[Suites décroissantes]
\explique{Pour montrer qu'une suite est décroissante, on montre que
$u_{n+1}-u_{n}\le 0$ pour tout $n$.}{propriété 1}
\cas{a)}
\begin{calculs}
u_{n+1} &= 7-2(n+1)^{2} &&\justif{on recopie l'expression}\\
u_{n+1} &= 7-2(n^{2}+2n+1) &&\justif{le carré d'abord, entre parenthèses}\\
u_{n+1} &= 7-2n^{2}-4n-2 &&\justif{le $-2$ multiplie chaque terme}\\
u_{n+1} &= -2n^{2}-4n+5 &&\\[4pt]
u_{n+1}-u_{n} &= -2n^{2}-4n+5-(7-2n^{2}) &&\justif{on soustrait $u_{n}$, entre parenthèses}\\
u_{n+1}-u_{n} &= -2n^{2}-4n+5-7+2n^{2} &&\\
u_{n+1}-u_{n} &= -4n-2 &&
\end{calculs}
\explique{Comme $n\ge 0$, on a $-4n-2\le -2<0$ : la suite $(u_{n})$ est
\res{strictement décroissante}.}{une différence négative pour tout $n$}
\cas{b)}
\begin{calculs}
v_{n+1}-v_{n} &= v_{n}-(n-3)^{2}-v_{n} &&\justif{on remplace $v_{n+1}$ par la relation}\\
v_{n+1}-v_{n} &= -(n-3)^{2} &&
\end{calculs}
\explique{Un carré est positif ou nul, donc son opposé est négatif ou nul :
$v_{n+1}-v_{n}\le 0$. La suite $(v_{n})$ est \res{décroissante}.}{pas besoin de
développer}
\explique{La différence s'annule pour $n=3$ : $v_{4}=v_{3}$. La suite n'est donc pas
strictement décroissante.}{$v_{3}=v_{4}=-9$}
\rem{« Décroissante » autorise deux termes égaux ; « strictement décroissante »
non.}
\cas{c)}
\begin{calculs}
w_{n+1} &= \dfrac{1}{(n+1)+1} &&\justif{on remplace $n$ par $(n+1)$}\\
w_{n+1} &= \dfrac{1}{n+2} &&\\[4pt]
w_{n+1}-w_{n} &= \dfrac{1}{n+2}-\dfrac{1}{n+1} &&\\
w_{n+1}-w_{n} &= \dfrac{1\fois{(n+1)}}{(n+2)\fois{(n+1)}}-\dfrac{1\fois{(n+2)}}{(n+1)\fois{(n+2)}}
  &&\justif{même dénominateur}\\
w_{n+1}-w_{n} &= \dfrac{(n+1)-(n+2)}{(n+1)(n+2)} &&\justif{parenthèses : le $-$ porte
sur $n+2$}\\
w_{n+1}-w_{n} &= \dfrac{-1}{(n+1)(n+2)} &&\justif{$n+1-n-2=-1$}
\end{calculs}
\explique{Le numérateur $-1$ est négatif, le dénominateur est un produit de deux nombres
positifs : la différence est négative. La suite $(w_{n})$ est
\res{strictement décroissante}.}{règle des signes d'un quotient}
\explique{On s'y attendait : plus le dénominateur est grand, plus la fraction
$\dfrac{1}{n+1}$ est petite.}{$w_{0}=1$ ; $w_{1}=0{,}5$ ; $w_{2}\approx 0{,}33$}
\end{exo}

\begin{exo}[Conjecturer, puis démontrer]
\cas{1.}
\begin{calculs}
u_{1} &= u_{0}+2\times 0-7 &&\justif{la relation avec $n=0$}\\
u_{1} &= 20+0-7 = 13 &&\\
u_{2} &= 13+2\times 1-7 = 8 &&\justif{$n=1$}\\
u_{3} &= 8+2\times 2-7 = 5 &&\justif{$n=2$}\\
u_{4} &= 5+2\times 3-7 = 4 &&\justif{$n=3$}\\
u_{5} &= 4+2\times 4-7 = 5 &&\justif{$n=4$}
\end{calculs}
\explique{$\res{u_{1}=13}$ ; $\res{u_{2}=8}$ ; $\res{u_{3}=5}$ ; $\res{u_{4}=4}$ ;
$\res{u_{5}=5}$. La suite semble décroître jusqu'à $u_{4}$, puis croître.}{une
conjecture, sur six termes}
\cas{2.}
\begin{calculs}
u_{n+1}-u_{n} &= u_{n}+2n-7-u_{n} &&\justif{on remplace $u_{n+1}$ par la relation}\\
u_{n+1}-u_{n} &= \res{2n-7} &&
\end{calculs}
\begin{calculs}
2n-7 &\ge 0 &&\justif{on cherche quand la différence est positive}\\
2n &\ge 7 &&\\
n &\ge 3{,}5 &&
\end{calculs}
\explique{$2n-7$ est négatif pour $n\le 3$ et positif pour $n\ge 4$.}{$n$ est entier}
\cas{3.}
\explique{Pour $n\ge 4$, $u_{n+1}-u_{n}>0$ : la suite est \res{strictement
croissante à partir du rang $4$}.}{avec $n=4$ : $u_{5}-u_{4}=1>0$}
\explique{La démonstration confirme la conjecture : $u_{4}=4$ est le plus petit
terme.}{avant le rang $4$, la suite décroît}
\end{exo}

\begin{exo}[Un carré toujours positif]
\cas{1.}
\begin{calculs}
v_{n+1}-v_{n} &= v_{n}+(v_{n}-1)^{2}-v_{n} &&\justif{on remplace $v_{n+1}$ par la
relation}\\
v_{n+1}-v_{n} &= (v_{n}-1)^{2} &&
\end{calculs}
\explique{Un carré est positif ou nul : pour tout $n$, $v_{n+1}-v_{n}\ge 0$. Donc
$\res{(v_{n}) \text{ est croissante}}$.}{inutile de connaître $v_{n}$ : le signe d'un
carré suffit}
\cas{2.}
\begin{calculs}
v_{1} &= v_{0}+(v_{0}-1)^{2} &&\justif{la relation avec $n=0$}\\
v_{1} &= 0+(0-1)^{2} = 1 &&\justif{$(-1)^{2}=1$}\\
v_{2} &= 1+(1-1)^{2} = 1 &&\justif{$n=1$ : $(1-1)^{2}=0$}\\
v_{3} &= 1+(1-1)^{2} = 1 &&\justif{$n=2$}
\end{calculs}
\explique{$\res{v_{1}=1}$ ; $\res{v_{2}=1}$ ; $\res{v_{3}=1}$ : dès que $v_{n}=1$, le carré $(v_{n}-1)^{2}$ est
nul et le terme suivant vaut encore $1$.}{la suite reste égale à $1$}
\explique{$v_{2}=v_{1}$ : la suite \res{n'est pas strictement croissante} ;
elle est croissante, et constante à partir du rang $1$.}{une égalité suffit à exclure
« strictement »}
\end{exo}

\partie{Sens de variation : quotient $\dfrac{u_{n+1}}{u_{n}}$}

\begin{exo}[Comparer le quotient à 1]
\explique{Si tous les termes sont strictement positifs, $u_{n+1}\ge u_{n}$ revient à
$\dfrac{u_{n+1}}{u_{n}}\ge 1$ : on vérifie d'abord le signe des termes, puis on
simplifie le quotient.}{propriété 2 ; commode quand $u_{n}$ contient une puissance}
\cas{a)}
\explique{$2>0$ et $5^{n}>0$, donc $u_{n}>0$ pour tout $n$.}{on vérifie d'abord le signe
des termes}
\begin{calculs}
\dfrac{u_{n+1}}{u_{n}} &= \dfrac{2\times 5^{n+1}}{2\times 5^{n}} &&\justif{on recopie}\\
\dfrac{u_{n+1}}{u_{n}} &= \dfrac{\barre{2}\times 5\times\barre{5^{n}}}{\barre{2}\times
\barre{5^{n}}} &&\justif{$5^{n+1}=5\times 5^{n}$}\\
\dfrac{u_{n+1}}{u_{n}} &= 5 &&
\end{calculs}
\explique{$5>1$ : la suite $(u_{n})$ est \res{strictement croissante}.}{chaque
terme est le quintuple du précédent}
\cas{b)}
\explique{$9>0$ et $0{,}8^{n}>0$, donc $v_{n}>0$ pour tout $n$.}{signe des termes}
\begin{calculs}
\dfrac{v_{n+1}}{v_{n}} &= \dfrac{9\times 0{,}8^{n+1}}{9\times 0{,}8^{n}}
  &&\justif{on recopie}\\
\dfrac{v_{n+1}}{v_{n}} &= \dfrac{\barre{9}\times 0{,}8\times\barre{0{,}8^{n}}}{\barre{9}
\times\barre{0{,}8^{n}}} &&\justif{$0{,}8^{n+1}=0{,}8\times 0{,}8^{n}$}\\
\dfrac{v_{n+1}}{v_{n}} &= 0{,}8 &&
\end{calculs}
\explique{$0{,}8<1$ : la suite $(v_{n})$ est \res{strictement
décroissante}.}{chaque terme vaut $80\,\%$ du précédent}
\rem{Le signe des termes compte : multiplier un nombre \emph{négatif} par $0{,}8$ le
rapproche de $0$, donc l'augmente. D'où la vérification $v_{n}>0$ avant de conclure.}
\cas{c)}
\explique{Pour $n\ge 1$, $4^{n}>0$ et $n>0$, donc $w_{n}>0$.}{signe des termes}
\begin{calculs}
\dfrac{w_{n+1}}{w_{n}} &= \dfrac{4^{n+1}}{n+1}\times\dfrac{n}{4^{n}} &&\justif{diviser
par $w_{n}$, c'est multiplier par son inverse}\\
\dfrac{w_{n+1}}{w_{n}} &= \dfrac{4\times\barre{4^{n}}\times n}{(n+1)\times\barre{4^{n}}}
  &&\justif{$4^{n+1}=4\times 4^{n}$}\\
\dfrac{w_{n+1}}{w_{n}} &= \dfrac{4n}{n+1} &&
\end{calculs}
\explique{Le quotient dépend de $n$ : on le compare à $1$ par une différence.}{signe de
$\frac{w_{n+1}}{w_{n}}-1$}
\begin{calculs}
\dfrac{w_{n+1}}{w_{n}}-1 &= \dfrac{4n}{n+1}-\dfrac{n+1}{n+1} &&\justif{$1=\frac{n+1}{n+1}$}\\
\dfrac{w_{n+1}}{w_{n}}-1 &= \dfrac{4n-(n+1)}{n+1} &&\justif{parenthèses : le $-$ porte
sur $n+1$}\\
\dfrac{w_{n+1}}{w_{n}}-1 &= \dfrac{3n-1}{n+1} &&
\end{calculs}
\explique{Pour $n\ge 1$ : $3n-1\ge 2>0$ et $n+1>0$, donc $\dfrac{w_{n+1}}{w_{n}}>1$. La
suite $(w_{n})$ est \res{strictement croissante}.}{$w_{1}=4$ ; $w_{2}=8$ ;
$w_{3}\approx 21{,}3$}
\end{exo}

\begin{exo}[D'autres quotients]
\cas{a)}
\explique{$3^{n}>0$ et $2^{n}>0$, donc $u_{n}>0$ pour tout $n$.}{signe des termes}
\begin{calculs}
\dfrac{u_{n+1}}{u_{n}} &= \dfrac{3^{n+1}}{2^{n+1}}\times\dfrac{2^{n}}{3^{n}}
  &&\justif{on multiplie par l'inverse de $u_{n}$}\\
\dfrac{u_{n+1}}{u_{n}} &= \dfrac{3\times\barre{3^{n}}\times\barre{2^{n}}}{2\times
\barre{2^{n}}\times\barre{3^{n}}} &&\justif{$3^{n+1}=3\times 3^{n}$ et
$2^{n+1}=2\times 2^{n}$}\\
\dfrac{u_{n+1}}{u_{n}} &= \dfrac{3}{2} = 1{,}5 &&
\end{calculs}
\explique{$1{,}5>1$ : la suite $(u_{n})$ est \res{strictement
croissante}.}{chaque terme est $1{,}5$ fois le précédent}
\cas{b)}
\explique{Pour $n\ge 1$, $n>0$ et $2^{n}>0$, donc $v_{n}>0$.}{signe des termes}
\begin{calculs}
\dfrac{v_{n+1}}{v_{n}} &= \dfrac{n+1}{2^{n+1}}\times\dfrac{2^{n}}{n} &&\justif{on
multiplie par l'inverse}\\
\dfrac{v_{n+1}}{v_{n}} &= \dfrac{(n+1)\times\barre{2^{n}}}{2\times\barre{2^{n}}\times n}
  &&\justif{$2^{n+1}=2\times 2^{n}$}\\
\dfrac{v_{n+1}}{v_{n}} &= \dfrac{n+1}{2n} &&\\[4pt]
\dfrac{v_{n+1}}{v_{n}}-1 &= \dfrac{n+1}{2n}-\dfrac{2n}{2n} &&\justif{$1=\frac{2n}{2n}$}\\
\dfrac{v_{n+1}}{v_{n}}-1 &= \dfrac{n+1-2n}{2n} &&\\
\dfrac{v_{n+1}}{v_{n}}-1 &= \dfrac{1-n}{2n} &&
\end{calculs}
\explique{Pour $n\ge 1$ : $1-n\le 0$ et $2n>0$, donc $\dfrac{v_{n+1}}{v_{n}}\le 1$. La
suite $(v_{n})$ est \res{décroissante}.}{quotient inférieur ou égal à $1$}
\explique{Le quotient vaut $1$ pour $n=1$ : $v_{1}=v_{2}=\dfrac{1}{2}$. La suite n'est
strictement décroissante qu'à partir du rang $2$.}{$v_{1}=\frac{1}{2}$ ;
$v_{2}=\frac{2}{4}$ ; $v_{3}=\frac{3}{8}$}
\cas{c)}
\explique{$5^{n}>0$ et $n+1>0$, donc $w_{n}>0$.}{signe des termes}
\begin{calculs}
w_{n+1} &= \dfrac{5^{n+1}}{(n+1)+1} = \dfrac{5^{n+1}}{n+2} &&\justif{on remplace $n$
par $(n+1)$ partout}\\[4pt]
\dfrac{w_{n+1}}{w_{n}} &= \dfrac{5^{n+1}}{n+2}\times\dfrac{n+1}{5^{n}} &&\justif{on
multiplie par l'inverse}\\
\dfrac{w_{n+1}}{w_{n}} &= \dfrac{5\times\barre{5^{n}}\times(n+1)}{(n+2)\times
\barre{5^{n}}} &&\justif{$5^{n+1}=5\times 5^{n}$}\\
\dfrac{w_{n+1}}{w_{n}} &= \dfrac{5n+5}{n+2} &&\\[4pt]
\dfrac{w_{n+1}}{w_{n}}-1 &= \dfrac{5n+5}{n+2}-\dfrac{n+2}{n+2} &&\justif{$1=\frac{n+2}{n+2}$}\\
\dfrac{w_{n+1}}{w_{n}}-1 &= \dfrac{5n+5-n-2}{n+2} &&\\
\dfrac{w_{n+1}}{w_{n}}-1 &= \dfrac{4n+3}{n+2} &&
\end{calculs}
\explique{$4n+3>0$ et $n+2>0$, donc $\dfrac{w_{n+1}}{w_{n}}>1$ : la suite $(w_{n})$ est
\res{strictement croissante}.}{quotient supérieur à $1$}
\rem{Erreur fréquente : écrire $w_{n+1}=\dfrac{5^{n+1}}{n+1}$, en oubliant de remplacer
le $n$ du dénominateur.}
\end{exo}

\begin{exo}[Choisir la méthode]
\explique{Produit avec une puissance $n$ et termes positifs : le quotient se simplifie.
Somme sans puissance : la différence.}{puissance $n$ : quotient (propriété 2) ; sinon : différence (propriété 1)}
\cas{a)}
\explique{Produit avec $0{,}5^{n}$ : le \textbf{quotient}. $6>0$ et $0{,}5^{n}>0$, donc
$u_{n}>0$.}{signe des termes}
\begin{calculs}
\dfrac{u_{n+1}}{u_{n}} &= \dfrac{6\times 0{,}5^{n+1}}{6\times 0{,}5^{n}} &&\justif{on
recopie}\\
\dfrac{u_{n+1}}{u_{n}} &= 0{,}5 &&\justif{$0{,}5^{n+1}=0{,}5\times 0{,}5^{n}$, on
simplifie}
\end{calculs}
\explique{$0{,}5<1$ et les termes sont positifs : la suite $(u_{n})$ est
\res{strictement décroissante}.}{chaque terme est la moitié du précédent}
\cas{b)}
\explique{Une somme, sans puissance $n$ : la \textbf{différence}.}{de plus $v_{0}=0$ :
on ne pourrait pas diviser par $v_{0}$}
\begin{calculs}
v_{n+1} &= (n+1)^{2}+5(n+1) &&\justif{on recopie l'expression}\\
v_{n+1} &= n^{2}+2n+1+5n+5 &&\justif{$(n+1)^{2}=n^{2}+2n+1$}\\
v_{n+1} &= n^{2}+7n+6 &&\\[4pt]
v_{n+1}-v_{n} &= n^{2}+7n+6-(n^{2}+5n) &&\justif{on soustrait $v_{n}$, entre parenthèses}\\
v_{n+1}-v_{n} &= 2n+6 &&\justif{les $n^{2}$ s'éliminent}
\end{calculs}
\explique{Comme $n\ge 0$, on a $2n+6>0$ : la suite $(v_{n})$ est
\res{strictement croissante}.}{une différence positive}
\cas{c)}
\explique{Produit avec $2^{n}$ : le \textbf{quotient}. Pour $n\ge 1$, $n>0$ et
$2^{n}>0$, donc $w_{n}>0$.}{signe des termes}
\begin{calculs}
\dfrac{w_{n+1}}{w_{n}} &= \dfrac{(n+1)\times 2^{n+1}}{n\times 2^{n}} &&\justif{on
recopie}\\
\dfrac{w_{n+1}}{w_{n}} &= \dfrac{(n+1)\times 2\times\barre{2^{n}}}{n\times\barre{2^{n}}}
  &&\justif{$2^{n+1}=2\times 2^{n}$}\\
\dfrac{w_{n+1}}{w_{n}} &= \dfrac{2n+2}{n} &&\\[4pt]
\dfrac{w_{n+1}}{w_{n}}-1 &= \dfrac{2n+2}{n}-\dfrac{n}{n} &&\justif{$1=\frac{n}{n}$}\\
\dfrac{w_{n+1}}{w_{n}}-1 &= \dfrac{n+2}{n} &&
\end{calculs}
\explique{$n+2>0$ et $n>0$ : le quotient est supérieur à $1$, la suite $(w_{n})$ est
\res{strictement croissante}.}{quotient supérieur à $1$}
\end{exo}

\begin{exo}[Croissante à partir d'un rang]
\cas{1.}
\begin{calculs}
u_{1} &= \dfrac{2^{1}}{1^{2}} = 2 &&\justif{on remplace $n$ par $1$}\\
u_{2} &= \dfrac{2^{2}}{2^{2}} = 1 &&\\
u_{3} &= \dfrac{2^{3}}{3^{2}} = \dfrac{8}{9} &&\\
u_{3} &\approx 0{,}89 &&\\
u_{4} &= \dfrac{2^{4}}{4^{2}} = \dfrac{16}{16} &&\\
u_{4} &= 1 &&\\
u_{5} &= \dfrac{2^{5}}{5^{2}} = \dfrac{32}{25} &&\\
u_{5} &= 1{,}28 &&
\end{calculs}
\explique{La suite décroît de $u_{1}$ à $u_{3}$, puis croît : on s'attend à un rang à
partir duquel elle est croissante.}{ce que la question 3 va préciser}
\cas{2.}
\explique{Pour $n\ge 1$, $2^{n}>0$ et $n^{2}>0$, donc $u_{n}>0$.}{signe des termes}
\begin{calculs}
\dfrac{u_{n+1}}{u_{n}} &= \dfrac{2^{n+1}}{(n+1)^{2}}\times\dfrac{n^{2}}{2^{n}}
  &&\justif{on multiplie par l'inverse}\\
\dfrac{u_{n+1}}{u_{n}} &= \dfrac{2\times\barre{2^{n}}\times n^{2}}{(n+1)^{2}\times
\barre{2^{n}}} &&\justif{$2^{n+1}=2\times 2^{n}$}\\
\dfrac{u_{n+1}}{u_{n}} &= \dfrac{2n^{2}}{(n+1)^{2}} &&
\end{calculs}
\explique{On transforme l'inégalité voulue étape par étape :}{chaque ligne équivaut à la
précédente}
\begin{calculs}
\dfrac{2n^{2}}{(n+1)^{2}} &\ge 1 &&\\
2n^{2} &\ge (n+1)^{2} &&\justif{on multiplie par $(n+1)^{2}>0$ : le sens ne change
pas}\\
2n^{2} &\ge n^{2}+2n+1 &&\justif{on développe}\\
n^{2}-2n-1 &\ge 0 &&\justif{on retire $n^{2}+2n+1$ aux deux membres}
\end{calculs}
\explique{$\dfrac{u_{n+1}}{u_{n}}\ge 1$ revient donc à $\res{n^{2}-2n-1\ge 0}$.}{ce qu'il
fallait montrer}
\cas{3.}
\explique{On étudie le signe du trinôme $x^{2}-2x-1$ : $a=1$, $b=-2$, $c=-1$.}{second
degré (chapitre 3)}
\begin{calculs}
\Delta &= b^{2}-4ac &&\\
\Delta &= (-2)^{2}-4\times 1\times(-1) = 8 &&\justif{$\Delta>0$ : deux racines}\\
x_{1} &= \dfrac{2-\sqrt{8}}{2} \approx -0{,}41 &&\justif{$x=\frac{-b\pm\sqrt{\Delta}}{2a}$}\\
x_{2} &= \dfrac{2+\sqrt{8}}{2} \approx 2{,}41 &&
\end{calculs}
\explique{$a=1>0$ : le trinôme est positif à l'extérieur des racines, donc pour
$x\ge x_{2}\approx 2{,}41$ quand $x$ est positif.}{signe de $a$ à l'extérieur des
racines}
\explique{Pour un entier $n\ge 1$ : $n^{2}-2n-1\ge 0$ dès $n\ge 3$, donc
$\dfrac{u_{n+1}}{u_{n}}\ge 1$, c'est-à-dire $u_{n+1}\ge u_{n}$.}{le premier entier après
$2{,}41$ est $3$}
\explique{La suite est \res{croissante à partir du rang $3$}.}{d'accord avec la
question 1 : $u_{3}<u_{4}<u_{5}$}
\end{exo}

\partie{Sens de variation : fonction associée}

\begin{exo}[Étudier la fonction associée]
\explique{Pour une suite explicite $u_{n}=f(n)$, les points $(n\,;u_{n})$ sont sur la
courbe de $f$ : si $f$ est croissante sur $\intfo{0}{+\infty}$, la suite l'est
aussi.}{propriété 3 : on étudie $f$ avec sa dérivée}
\cas{a)}
\explique{On pose $f(x)=x^{3}+6x^{2}+1$ sur $\intfo{0}{+\infty}$ : $u_{n}=f(n)$.}{la
fonction associée : on remplace $n$ par $x$}
\begin{calculs}
f'(x) &= 3x^{2}+12x &&\justif{$(x^{3})'=3x^{2}$ ; $(x^{2})'=2x$ ; $(1)'=0$}\\
f'(x) &= 3x(x+4) &&\justif{on factorise par $3x$ pour le signe}
\end{calculs}
\explique{Pour $x\ge 0$ : $3x\ge 0$ et $x+4>0$, donc $f'(x)\ge 0$ ; $f'$ ne s'annule
qu'en $0$, donc $f$ est strictement croissante sur $\intfo{0}{+\infty}$.}{signe d'un
produit}
\explique{La suite $(u_{n})$ est donc \res{strictement croissante}.}{propriété 3}
\cas{b)}
\explique{On pose $f(x)=\dfrac{2x+3}{x+1}$ sur $\intfo{0}{+\infty}$.}{la fonction
associée}
\begin{calculs}
f'(x) &= \dfrac{2(x+1)-(2x+3)\times 1}{(x+1)^{2}} &&\justif{$\left(\frac{u}{v}\right)'
=\frac{u'v-uv'}{v^{2}}$, $u=2x+3$, $v=x+1$}\\
f'(x) &= \dfrac{2x+2-2x-3}{(x+1)^{2}} &&\justif{parenthèses : le $-$ porte sur $2x+3$}\\
f'(x) &= \dfrac{-1}{(x+1)^{2}} &&
\end{calculs}
\explique{$-1<0$ et $(x+1)^{2}>0$ : $f'(x)<0$, donc $f$ est strictement décroissante sur
$\intfo{0}{+\infty}$.}{un carré non nul est positif}
\explique{La suite $(v_{n})$ est donc \res{strictement décroissante}.}{$v_{0}=3$ ;
$v_{1}=2{,}5$ ; $v_{2}\approx 2{,}33$}
\cas{c)}
\explique{On pose $f(x)=x^{3}-3x^{2}$ sur $\intfo{0}{+\infty}$.}{la fonction associée}
\begin{calculs}
f'(x) &= 3x^{2}-6x &&\justif{on dérive}\\
f'(x) &= 3x(x-2) &&\justif{on factorise par $3x$}\\
f(0) &= 0 &&\\
f(2) &= 8-12 = -4 &&\justif{image de la racine $2$}
\end{calculs}
\explique{Pour $x\ge 0$, on a $3x\ge 0$ : $f'(x)$ a le signe de $x-2$, négatif avant $2$,
positif après.}{on lit le signe facteur par facteur}
\begin{center}
\begin{tikzpicture}[scale=0.8]
  \tkzTabInit[lgt=2,espcl=2.4]{$x$/0.8,$f'(x)$/0.8,$f(x)$/1.4}{$0$,$2$,$+\infty$}
  \tkzTabLine{z,-,z,+,}
  \tkzTabVar{+/$0$,-/$-4$,+/{}}
\end{tikzpicture}
\end{center}
\explique{$f$ est strictement croissante sur $\intfo{2}{+\infty}$ : la suite $(w_{n})$ est
\res{strictement croissante à partir du rang $2$}.}{avant, elle décroît :
$w_{0}=0$ ; $w_{1}=-2$ ; $w_{2}=-4$ ; $w_{3}=0$}
\end{exo}

\begin{exo}[D'autres fonctions]
\cas{a)}
\explique{On pose $f(x)=x^{2}-8x+3$ sur $\intfo{0}{+\infty}$.}{la fonction associée}
\begin{calculs}
f'(x) &= 2x-8 &&\justif{on dérive}\\
f(4) &= 16-32+3 = -13 &&\justif{$f'(x)=0$ pour $x=4$}
\end{calculs}
\explique{$2x-8\ge 0$ revient à $x\ge 4$ : $f'$ est négative sur $\intff{0}{4}$, positive
sur $\intfo{4}{+\infty}$.}{inéquation du premier degré}
\begin{center}
\begin{tikzpicture}[scale=0.8]
  \tkzTabInit[lgt=2,espcl=2.4]{$x$/0.8,$f'(x)$/0.8,$f(x)$/1.4}{$0$,$4$,$+\infty$}
  \tkzTabLine{,-,z,+,}
  \tkzTabVar{+/$3$,-/$-13$,+/{}}
\end{tikzpicture}
\end{center}
\explique{La suite $(u_{n})$ est \res{strictement croissante à partir du rang $4$} ; avant, elle décroît.}{$u_{3}=-12$ ; $u_{4}=-13$ ; $u_{5}=-12$}
\cas{b)}
\explique{On pose $f(x)=5-\dfrac{4}{x+2}$ sur $\intfo{0}{+\infty}$.}{la fonction
associée}
\begin{calculs}
f'(x) &= 0-4\times\left(-\dfrac{1}{(x+2)^{2}}\right) &&\justif{$(5)'=0$ ;
$\left(\frac{1}{v}\right)'=-\frac{v'}{v^{2}}$ avec $v=x+2$, $v'=1$}\\
f'(x) &= \dfrac{4}{(x+2)^{2}} &&\justif{moins par moins : plus}
\end{calculs}
\explique{$4>0$ et $(x+2)^{2}>0$ : $f'(x)>0$, $f$ est strictement croissante. La suite
$(v_{n})$ est \res{strictement croissante}.}{propriété 3}
\rem{Sans dérivée : $\dfrac{4}{n+2}$ diminue quand $n$ augmente, donc on retire de $5$ un
nombre de plus en plus petit, et $v_{n}$ augmente.}
\cas{c)}
\explique{On pose $f(x)=\dfrac{x}{x^{2}+1}$ sur $\intfo{0}{+\infty}$ ; $u=x$, $u'=1$,
$v=x^{2}+1$, $v'=2x$.}{dérivée d'un quotient}
\begin{calculs}
f'(x) &= \dfrac{1\times(x^{2}+1)-x\times 2x}{(x^{2}+1)^{2}}
  &&\justif{$\frac{u'v-uv'}{v^{2}}$}\\
f'(x) &= \dfrac{x^{2}+1-2x^{2}}{(x^{2}+1)^{2}} &&\\
f'(x) &= \dfrac{1-x^{2}}{(x^{2}+1)^{2}} &&\\
f'(x) &= \dfrac{(1-x)(1+x)}{(x^{2}+1)^{2}} &&\justif{$a^{2}-b^{2}=(a-b)(a+b)$}\\
f(1) &= \dfrac{1}{1+1} = 0{,}5 &&\justif{image de la racine $1$}
\end{calculs}
\explique{Pour $x\ge 0$ : $1+x>0$ et le dénominateur est positif ; $f'(x)$ a le signe de
$1-x$, positif avant $1$, négatif après.}{un seul facteur change de signe}
\begin{center}
\begin{tikzpicture}[scale=0.8]
  \tkzTabInit[lgt=2,espcl=2.4]{$x$/0.8,$f'(x)$/0.8,$f(x)$/1.4}{$0$,$1$,$+\infty$}
  \tkzTabLine{,+,z,-,}
  \tkzTabVar{-/$0$,+/$0{,}5$,-/{}}
\end{tikzpicture}
\end{center}
\explique{La suite $(w_{n})$ est \res{strictement décroissante à partir du rang $1$}.}{$w_{0}=0$ ; $w_{1}=0{,}5$ ; $w_{2}=0{,}4$ ; $w_{3}=0{,}3$}
\end{exo}

\begin{exo}[Une erreur fréquente]
\cas{1.}
\begin{calculs}
u_{1} &= f(u_{0}) = 0{,}5\times 10+2 &&\justif{$u_{n+1}=f(u_{n})$ avec $n=0$}\\
u_{1} &= 5+2 = \res{7} &&\\
u_{2} &= 0{,}5\times 7+2 = \res{5{,}5} &&\justif{$u_{2}=f(u_{1})$}\\
u_{3} &= 0{,}5\times 5{,}5+2 = \res{4{,}75} &&\justif{$u_{3}=f(u_{2})$}
\end{calculs}
\cas{2.}
\explique{$u_{0}>u_{1}>u_{2}>u_{3}$ : la suite décroît, alors que $f$ est croissante.
\res{L'affirmation est fausse}.}{un contre-exemple suffit}
\explique{La propriété 3 ne vaut que pour une suite explicite $u_{n}=f(n)$ : les points
$(n\,;f(n))$ sont alors sur la courbe de $f$.}{ce que dit le cours}
\explique{Ici $u_{n+1}=f(u_{n})$ : la croissance de $f$ ne dit rien de la comparaison
entre $u_{n+1}$ et $u_{n}$.}{suite définie par récurrence}
\rem{Pour une suite définie par récurrence, on revient à la différence
$u_{n+1}-u_{n}$ (partie VII).}
\end{exo}

\begin{exo}[Suite croissante, fonction non croissante]
\cas{1.}
\begin{calculs}
f(x) &= x^{2}-0{,}8x+0{,}16 &&\justif{on développe $(x-0{,}4)^{2}$}\\
f'(x) &= 2x-0{,}8 &&\justif{on dérive}\\
f'(x) &= 2(x-0{,}4) &&\justif{on factorise par $2$}\\
f(0) &= (0-0{,}4)^{2} = 0{,}16 &&\\
f(0{,}4) &= 0 &&
\end{calculs}
\explique{$f'(x)$ a le signe de $x-0{,}4$ : négatif sur $\intff{0}{0{,}4}$, positif
ensuite.}{$2>0$}
\begin{center}
\begin{tikzpicture}[scale=0.8]
  \tkzTabInit[lgt=2,espcl=2.4]{$x$/0.8,$f'(x)$/0.8,$f(x)$/1.4}{$0$,$0{,}4$,$+\infty$}
  \tkzTabLine{,-,z,+,}
  \tkzTabVar{+/$0{,}16$,-/$0$,+/{}}
\end{tikzpicture}
\end{center}
\explique{$f$ est décroissante sur $\intff{0}{0{,}4}$ et croissante sur
$\intfo{0{,}4}{+\infty}$.}{lecture du tableau}
\cas{2.}
\begin{calculs}
u_{n+1} &= (n+1-0{,}4)^{2} &&\justif{on remplace $n$ par $(n+1)$}\\
u_{n+1} &= (n+0{,}6)^{2} &&\\
u_{n+1} &= n^{2}+1{,}2n+0{,}36 &&\justif{$(a+b)^{2}$ : $2\times 0{,}6=1{,}2$ ;
$0{,}6^{2}=0{,}36$}\\[4pt]
u_{n+1}-u_{n} &= n^{2}+1{,}2n+0{,}36-(n^{2}-0{,}8n+0{,}16) &&\justif{$u_{n}$ développé
à la question 1}\\
u_{n+1}-u_{n} &= n^{2}+1{,}2n+0{,}36-n^{2}+0{,}8n-0{,}16 &&\\
u_{n+1}-u_{n} &= 2n+0{,}2 &&
\end{calculs}
\explique{Comme $n\ge 0$, on a $2n+0{,}2>0$ : la suite $(u_{n})$ est
\res{strictement croissante}.}{$u_{0}=0{,}16$ ; $u_{1}=0{,}36$ ;
$u_{2}=2{,}56$}
\cas{3.}
\explique{\res{Non} : ici la suite est croissante alors que $f$ n'est pas
croissante sur $\intfo{0}{+\infty}$ (elle décroît sur $\intff{0}{0{,}4}$).}{un
contre-exemple}
\explique{La propriété 3 ne marche que dans un sens : « $f$ croissante » entraîne
« suite croissante », mais pas l'inverse.}{la réciproque est fausse}
\rem{La suite ne « voit » que les entiers $0$ ; $1$ ; $2$\dots{} La baisse de $f$ entre $0$
et $0{,}4$ se passe entre deux entiers : elle est invisible pour la suite.}
\end{exo}

\partie{Notion de limite}

\begin{exo}[Conjecturer une limite]
\explique{On calcule des termes de rang de plus en plus grand et on regarde de quoi ils
se rapprochent.}{une conjecture, pas une démonstration}
\cas{a)}
\begin{calculs}
u_{9} &= 2+\dfrac{5}{9+1} = 2{,}5 &&\justif{on remplace $n$ par $9$ ; $\frac{5}{10}=0{,}5$}\\
u_{99} &= 2+\dfrac{5}{100} = 2{,}05 &&\justif{$n=99$}\\
u_{999} &= 2+\dfrac{5}{1\,000} = 2{,}005 &&\justif{$n=999$}
\end{calculs}
\explique{Les termes se rapprochent de $2$ : on conjecture
$\displaystyle\lim_{n\to+\infty}u_{n}=\res{2}$.}{$\frac{5}{n+1}$ devient très petit}
\cas{b)}
\begin{calculs}
v_{10} &= 3\times 10^{2}+1 = 301 &&\justif{on remplace $n$ par $10$}\\
v_{100} &= 3\times 100^{2}+1 = 30\,001 &&\\
v_{1\,000} &= 3\times 1\,000^{2}+1 = 3\,000\,001 &&
\end{calculs}
\explique{Les termes deviennent aussi grands que l'on veut : on conjecture
$\displaystyle\lim_{n\to+\infty}v_{n}=\res{+\infty}$.}{$n^{2}$ grandit sans fin}
\cas{c)}
\explique{$w_{0}=4$ ; $w_{1}=-4$ ; $w_{2}=4$ ; $w_{3}=-4$.}{$(-1)^{n}$ vaut $1$ si $n$ est
pair, $-1$ sinon}
\explique{Les termes alternent entre $4$ et $-4$ : la suite $(w_{n})$
\res{n'a pas de limite}.}{ils ne se rapprochent d'aucun nombre}
\rem{« Pas de limite » ne veut dire ni « limite $0$ » ni « limite $4$ ou $-4$ » : une
limite est un seul nombre, dont les termes deviennent aussi proches que l'on veut.}
\end{exo}

\begin{exo}[Limite finie ou infinie]
\cas{a)}
\begin{calculs}
a_{10} &= 5-10^{2} = -95 &&\justif{on remplace $n$ par $10$}\\
a_{100} &= 5-100^{2} = -9\,995 &&\\
a_{1\,000} &= 5-1\,000^{2} = -999\,995 &&
\end{calculs}
\explique{Les termes sont négatifs et s'éloignent de $0$ autant que l'on veut : on
conjecture $\displaystyle\lim_{n\to+\infty}a_{n}=\res{-\infty}$.}{$-n^{2}$ descend sans
fin}
\cas{b)}
\begin{calculs}
b_{10} &= \dfrac{1}{10^{2}} = 0{,}01 &&\justif{$10^{2}=100$}\\
b_{100} &= \dfrac{1}{100^{2}} = 0{,}000\,1 &&\justif{$100^{2}=10\,000$}\\
b_{1\,000} &= \dfrac{1}{1\,000^{2}} = 0{,}000\,001 &&\justif{$1\,000^{2}=1\,000\,000$}
\end{calculs}
\explique{Les termes, positifs, se rapprochent de $0$ : on conjecture
$\displaystyle\lim_{n\to+\infty}b_{n}=\res{0}$.}{$1$ divisé par un nombre de plus en
plus grand}
\cas{c)}
\begin{calculs}
c_{10} &= \dfrac{2\times 10+1}{10} = 2{,}1 &&\justif{$\frac{21}{10}$}\\
c_{100} &= \dfrac{201}{100} = 2{,}01 &&\\
c_{1\,000} &= \dfrac{2\,001}{1\,000} = 2{,}001 &&
\end{calculs}
\explique{Les termes se rapprochent de $2$ : on conjecture
$\displaystyle\lim_{n\to+\infty}c_{n}=\res{2}$.}{$c_{n}=\frac{2n}{n}+\frac{1}{n}=2+\frac{1}{n}$}
\end{exo}

\begin{exo}[Avec la calculatrice]
\cas{1.}
\explique{\texttt{Récurrente d'ordre 1} : $u_{n+1}=0{,}5\,u_{n}+3$ et $u_{0}=1$, puis
\texttt{Tableau} jusqu'à $30$.}{contrôle à la main : $u_{1}=0{,}5+3=3{,}5$}
\explique{$\res{u_{10}\approx 5{,}995\,117}$ ; $\res{u_{20}\approx 5{,}999\,995}$ ;
$\res{u_{30}\approx 6{,}000\,000}$.}{arrondis à $10^{-6}$}
\cas{2.}
\explique{Les termes se rapprochent de $6$ : on conjecture
$\res{\displaystyle\lim_{n\to+\infty}u_{n}=6}$.}{l'écart à $6$ devient aussi petit que
l'on veut}
\rem{$u_{30}$ n'est pas égal à $6$ : c'est son arrondi à $10^{-6}$ qui vaut $6{,}000\,000$.
L'écart à $6$ est divisé par $2$ à chaque rang, sans jamais s'annuler.}
\end{exo}

\begin{exo}[Alterner et converger]
\cas{1.}
\begin{calculs}
u_{10} &= \dfrac{(-1)^{10}}{10} = \dfrac{1}{10} &&\justif{$10$ est pair :
$(-1)^{10}=1$}\\
u_{10} &= \res{0{,}1} &&\\
u_{11} &= \dfrac{(-1)^{11}}{11} = -\dfrac{1}{11} &&\justif{$11$ est impair :
$(-1)^{11}=-1$}\\
u_{11} &\approx \res{-0{,}091} &&\\
u_{100} &= \dfrac{1}{100} = \res{0{,}01} &&\justif{$100$ est pair}\\
u_{101} &= -\dfrac{1}{101} \approx \res{-0{,}009\,9} &&\justif{$101$ est impair}
\end{calculs}
\cas{2.}
\explique{Oui, la suite change de signe à chaque rang : positive pour $n$ pair, négative
pour $n$ impair.}{le signe vient de $(-1)^{n}$}
\explique{Mais les termes s'approchent de $0$ des deux côtés, de plus en plus près : on
conjecture $\res{\displaystyle\lim_{n\to+\infty}u_{n}=0}$.}{la distance à $0$ est
$\frac{1}{n}$, qui devient très petite}
\rem{À comparer avec l'exercice 37 c) : là, les termes alternent entre $4$ et $-4$ sans
se rapprocher de rien. Une suite qui alterne peut avoir une limite, à condition que les
écarts à ce nombre deviennent aussi petits que l'on veut.}
\end{exo}

\partie{Problèmes : déterminer un seuil}

\begin{exo}[Abonnés d'une newsletter]
\cas{1.}
\explique{$u_{0}=3\,000$.}{en janvier : $n=0$}
\begin{calculs}
u_{1} &= 3\,000-0{,}15\times 3\,000+800 &&\justif{on retire $15\,\%$, on ajoute les
inscrits}\\
u_{1} &= 3\,000-450+800 = \res{3\,350} &&\justif{$15\,\%$ de $3\,000$ : $450$}\\[4pt]
u_{2} &= 3\,350-0{,}15\times 3\,350+800 &&\justif{même calcul à partir de $u_{1}$}\\
u_{2} &= 3\,350-502{,}5+800 = \res{3\,647{,}5} &&
\end{calculs}
\rem{Un nombre d'abonnés est entier : $u_{2}$ donne environ $3\,648$ abonnés. La suite
est un modèle.}
\cas{2.}
\explique{Si $15\,\%$ des abonnés partent, il en reste $85\,\%$ : on multiplie par
$1-0{,}15=0{,}85$. Puis on ajoute les $800$ inscrits :
$\res{u_{n+1}=0{,}85\,u_{n}+800}$.}{coefficient multiplicateur, puis ajout}
\explique{Vérification : $0{,}85\times 3\,000+800=2\,550+800=3\,350$, c'est bien
$u_{1}$.}{la relation redonne la question 1}
\cas{3.}
\begin{minipage}[t]{0.42\linewidth}\vspace{0pt}
\begin{python}
def seuil():
    u = 3000
    n = 0
    while u <= 5000 :
        u = 0.85*u + 800
        n = n + 1
    return n
\end{python}
\end{minipage}\hfill
\begin{minipage}[t]{0.54\linewidth}\vspace{0pt}\raggedright
\gris{\texttt{u} contient le terme, \texttt{n} son rang ; tant que le seuil n'est pas
dépassé ($u\le 5\,000$, le contraire de « dépasser »), on passe au terme suivant
(relation de récurrence) et au mois suivant ; la boucle s'arrête au premier terme qui
dépasse $5\,000$, et \texttt{n} est son rang}
\end{minipage}
\rem{Erreur fréquente : écrire \texttt{while u > 5000}. Au départ $u=3\,000$, la boucle
ne tourne pas et la fonction renvoie $0$. La condition du \texttt{while} est le
contraire du but.}
\cas{4.}
\explique{\texttt{Tableau} de la calculatrice : $u_{11}\approx 4\,942{,}9$ et
$u_{12}\approx 5\,001{,}4$.}{le premier terme au-dessus de $5\,000$ est $u_{12}$}
\explique{\texttt{seuil()} renvoie $12$ : la newsletter dépasse $5\,000$ abonnés
\res{au bout de $12$ mois}, en janvier de l'année suivante.}{on conclut
par une phrase, dans le contexte}
\end{exo}

\begin{exo}[Seuil d'une suite explicite]
\cas{1.}
\begin{minipage}[t]{0.42\linewidth}\vspace{0pt}
\begin{python}
def seuil():
    n = 0
    u = 0
    while u <= 1000 :
        n = n + 1
        u = 2*n**2 + 3*n
    return n
\end{python}
\end{minipage}\hfill
\begin{minipage}[t]{0.54\linewidth}\vspace{0pt}\raggedright
\gris{on veut $u_{n}>1\,000$ : on continue tant que $u\le 1\,000$ ; la suite est
explicite : on calcule le terme directement à partir du rang, \texttt{u = 2*n**2 + 3*n}
(\texttt{**} est la puissance), sans utiliser le terme précédent}
\end{minipage}
\explique{L'ordre compte : \texttt{n} augmente d'abord, puis \texttt{u} est calculé avec
le nouveau rang ; \texttt{u} contient donc toujours $u_{n}$.}{au départ $u_{0}=0$}
\cas{2.}
\explique{\texttt{Explicite} : $u_{n}=2n^{2}+3n$, puis \texttt{Tableau}.}{on cherche le
premier terme au-dessus de $1\,000$}
\begin{calculs}
u_{21} &= 2\times 21^{2}+3\times 21 = 945 &&\justif{$2\times 441+63$ ; encore $\le 1\,000$}\\
u_{22} &= 2\times 22^{2}+3\times 22 = 1\,034 &&\justif{$2\times 484+66$ ; au-dessus de
$1\,000$}
\end{calculs}
\explique{Le plus petit rang tel que $u_{n}>1\,000$ est $\res{n=22}$.}{la fonction
renvoie $22$}
\end{exo}

\begin{exo}[Élimination d'un médicament]
\cas{1.}
\begin{calculs}
m_{1} &= 40-0{,}2\times 40 &&\justif{on élimine $20\,\%$ des $40$~mg}\\
m_{1} &= 40-8 = \res{32} &&\\
m_{2} &= 32-0{,}2\times 32 &&\\
m_{2} &= 32-6{,}4 = \res{25{,}6} &&
\end{calculs}
\explique{Éliminer $20\,\%$, c'est multiplier par $1-0{,}2=0{,}8$ :
$\res{m_{n+1}=0{,}8\,m_{n}}$.}{vérification : $0{,}8\times 40=32$}
\cas{2.}
\begin{minipage}[t]{0.42\linewidth}\vspace{0pt}
\begin{python}
def seuil():
    m = 40
    n = 0
    while m >= 5 :
        m = 0.8*m
        n = n + 1
    return n
\end{python}
\end{minipage}\hfill
\begin{minipage}[t]{0.54\linewidth}\vspace{0pt}\raggedright
\gris{on veut « moins de $5$~mg », soit $m<5$ ; on continue donc tant que $m\ge 5$ ; à
chaque tour, une heure passe : la quantité est multipliée par $0{,}8$ et le compteur
avance}
\end{minipage}
\cas{3.}
\explique{\texttt{Tableau} : $m_{9}\approx 5{,}37$ et $m_{10}\approx 4{,}29$.}{le premier
terme sous $5$ est $m_{10}$}
\explique{\texttt{seuil()} renvoie $10$ : il reste moins de $5$~mg
\res{au bout de $10$ heures}.}{on conclut dans le contexte}
\end{exo}

\begin{exo}[Qui dépasse l'autre ?]
\cas{1.}
\explique{\texttt{b} contient $b_{n}$ et \texttt{n} le mois ; \texttt{500 + 40*n} est
$a_{n}$, calculé directement : la suite $(a_{n})$ est explicite, elle n'a pas besoin de
variable.}{deux suites, une seule variable}
\explique{Tant que $b_{n}\le a_{n}$, on passe au mois suivant : \texttt{b} est multiplié
par $1{,}1$ et \texttt{n} augmente de $1$.}{la condition du \texttt{while}}
\explique{La fonction renvoie le \res{premier mois où le second site a plus
d'abonnés que le premier}.}{la boucle s'arrête dès que $b_{n}>a_{n}$}
\cas{2.}
\explique{\texttt{Tableau} avec les deux suites :}{on compare ligne par ligne}
\begin{calculs}
a_{19} &= 500+40\times 19 = 1\,260 &&\justif{et $b_{19}\approx 1\,223{,}2$ : encore
$b_{19}\le a_{19}$}\\
a_{20} &= 500+40\times 20 = 1\,300 &&\justif{et $b_{20}\approx 1\,345{,}5$ : cette fois
$b_{20}>a_{20}$}
\end{calculs}
\explique{La fonction renvoie $\res{20}$ : au bout de $20$ mois, le site qui gagne
$10\,\%$ par mois compte plus d'abonnés que celui qui en gagne $40$.}{interprétation}
\rem{Le premier site gagne toujours $40$ abonnés ; le second gagne $10\,\%$ d'un nombre
qui grandit : son gain mensuel augmente, il finit par rattraper l'autre.}
\end{exo}

\partiesansnum{Exercices type bac}

\begin{exo}[Automatismes (QCM)]
\cas{1.}
\begin{calculs}
u_{n+1} &= (n+1)^{2}-1 &&\justif{on remplace $n$ par $(n+1)$}\\
u_{n+1} &= n^{2}+2n+1-1 &&\\
u_{n+1} &= n^{2}+2n &&
\end{calculs}
\explique{Réponse \res{b)}.}{la réponse a) est $u_{n}+1=n^{2}$ : le piège}
\cas{2.}
\begin{calculs}
u_{1} &= 3\times 2-1 = 5 &&\justif{$n=0$}\\
u_{2} &= 3\times 5-1 = 14 &&\justif{$n=1$}
\end{calculs}
\explique{Réponse \res{b)}.}{a) est $u_{1}$, pas $u_{2}$}
\cas{3.}
\explique{Baisser de $15\,\%$, c'est multiplier par $1-0{,}15=0{,}85$ : réponse
\res{c)}.}{$1{,}15$ correspondrait à une hausse}
\cas{4.}
\explique{$-n^{2}-1=-(n^{2}+1)<0$ pour tout $n$ : $u_{n+1}<u_{n}$, la suite est
décroissante. Réponse \res{b)}.}{signe de la différence}
\cas{5.}
\explique{$\dfrac{1}{n}$ devient très petit, donc $u_{n}$ se rapproche de $3$ : réponse
\res{b)}.}{par exemple $u_{1\,000}=3{,}001$}
\end{exo}

\begin{exo}[Une plateforme de covoiturage]
\cas{1.}
\begin{calculs}
u_{1} &= 600-0{,}1\times 600+40 &&\justif{$10\,\%$ partent, $40$ arrivent}\\
u_{1} &= 600-60+40 = \res{580} &&\\
u_{2} &= 580-58+40 = \res{562} &&\justif{$10\,\%$ de $580$ : $58$}
\end{calculs}
\explique{Si $10\,\%$ des conducteurs partent, il en reste $90\,\%$ : on multiplie par
$0{,}9$, puis on ajoute $40$. Donc $u_{n+1}=0{,}9\,u_{n}+40$.}{coefficient
multiplicateur, puis ajout}
\cas{2.}
\begin{calculs}
400+200\times 0{,}9^{0} &= 400+200 = 600 &&\justif{$n=0$ : $0{,}9^{0}=1$}\\
400+200\times 0{,}9^{1} &= 400+180 = 580 &&\justif{$n=1$ : $200\times 0{,}9=180$}
\end{calculs}
\explique{On retrouve $u_{0}=600$ et $u_{1}=580$ : la formule est vérifiée pour $n=0$ et
$n=1$.}{on compare avec la question 1}
\cas{3.}
\begin{calculs}
u_{n+1} &= 400+200\times 0{,}9^{n+1} &&\justif{on remplace $n$ par $n+1$ dans la
formule admise}\\[4pt]
u_{n+1}-u_{n} &= 400+200\times 0{,}9^{n+1}-(400+200\times 0{,}9^{n}) &&\\
u_{n+1}-u_{n} &= 200\times 0{,}9^{n+1}-200\times 0{,}9^{n} &&\justif{les $400$
s'éliminent}\\
u_{n+1}-u_{n} &= 200\times 0{,}9\times 0{,}9^{n}-200\times 0{,}9^{n}
  &&\justif{$0{,}9^{n+1}=0{,}9\times 0{,}9^{n}$}\\
u_{n+1}-u_{n} &= 180\times 0{,}9^{n}-200\times 0{,}9^{n} &&\\
u_{n+1}-u_{n} &= -20\times 0{,}9^{n} &&\justif{$180-200=-20$}
\end{calculs}
\explique{$0{,}9^{n}>0$ (puissance d'un nombre positif), donc $-20\times 0{,}9^{n}<0$ :
la suite $(u_{n})$ est \res{strictement décroissante}.}{une différence négative}
\explique{Le nombre de conducteurs diminue chaque mois.}{interprétation}
\cas{4.}
\explique{Quand $n$ grandit, $0{,}9^{n}$ devient très petit : $0{,}9^{10}\approx 0{,}35$ ;
$0{,}9^{50}\approx 0{,}005$ ; $0{,}9^{100}\approx 0{,}000\,03$.}{on multiplie sans cesse
par un nombre inférieur à $1$}
\explique{$200\times 0{,}9^{n}$ se rapproche de $0$ : on conjecture
$\res{\displaystyle\lim_{n\to+\infty}u_{n}=400}$.}{$u_{n}=400+\text{un nombre de plus en
plus petit}$}
\explique{À long terme, le nombre de conducteurs diminue et se stabilise autour de
$400$, sans descendre en dessous.}{interprétation}
\cas{5.}
\begin{minipage}[t]{0.42\linewidth}\vspace{0pt}
\begin{python}
def seuil():
    u = 600
    n = 0
    while u >= 450 :
        u = 0.9*u + 40
        n = n + 1
    return n
\end{python}
\end{minipage}\hfill
\begin{minipage}[t]{0.54\linewidth}\vspace{0pt}\raggedright
\gris{on veut « moins de $450$ » ($u<450$) : on continue tant que $u\ge 450$ ; la ligne
de la boucle est la relation de récurrence de la question 1}
\end{minipage}
\explique{\texttt{Tableau} : $u_{13}\approx 450{,}8$ et $u_{14}\approx 445{,}8$ ; la
fonction renvoie $14$.}{le premier terme sous $450$ est $u_{14}$}
\explique{$n=12$ correspond à janvier 2027, donc $n=14$ à mars 2027 : il y a moins de
$450$ conducteurs pour la première fois en \res{mars 2027}.}{janvier 2026 :
$n=0$}
\end{exo}

\begin{exo}[Coût moyen de fabrication]
\cas{1.}
\begin{calculs}
c_{1} &= 1+\dfrac{400}{1} = \res{401} &&\justif{on remplace $n$ par $1$}\\
c_{10} &= 10+\dfrac{400}{10} = \res{50} &&\justif{$10+40$}\\
c_{20} &= 20+\dfrac{400}{20} = \res{40} &&\justif{$20+20$}
\end{calculs}
\explique{Pour $1$ lampe par jour, chaque lampe coûte $401$~€ ; pour $20$ lampes,
$40$~€.}{interprétation}
\cas{2.}
\begin{calculs}
f'(x) &= 1-\dfrac{400}{x^{2}} &&\justif{$(x)'=1$ ; $\left(\frac{1}{x}\right)'=-\frac{1}{x^{2}}$,
donc $\left(\frac{400}{x}\right)'=-\frac{400}{x^{2}}$}\\
f'(x) &= \dfrac{x^{2}}{x^{2}}-\dfrac{400}{x^{2}} &&\justif{$1=\frac{x^{2}}{x^{2}}$ :
même dénominateur}\\
f'(x) &= \dfrac{x^{2}-400}{x^{2}} &&\\
f'(x) &= \dfrac{(x-20)(x+20)}{x^{2}} &&\justif{$a^{2}-b^{2}=(a-b)(a+b)$ avec
$b=20$}
\end{calculs}
\explique{Sur $\intfo{1}{+\infty}$ : $x+20>0$ et $x^{2}>0$, donc $f'(x)$ a le signe de
$x-20$.}{un seul facteur change de signe}
\begin{center}
\begin{tikzpicture}[scale=0.8]
  \tkzTabInit[lgt=2,espcl=2.4]{$x$/0.8,$f'(x)$/0.8,$f(x)$/1.4}{$1$,$20$,$+\infty$}
  \tkzTabLine{,-,z,+,}
  \tkzTabVar{+/$401$,-/$40$,+/{}}
\end{tikzpicture}
\end{center}
\explique{$f(1)=401$ et $f(20)=40$ : ce sont $c_{1}$ et $c_{20}$.}{valeurs de la
question 1}
\cas{3.}
\explique{$c_{n}=f(n)$ : $(c_{n})$ est décroissante jusqu'au rang $20$, puis
\res{croissante à partir du rang $20$}.}{propriété 3, sur chaque intervalle}
\explique{Le plus petit terme est $c_{20}=40$ : le coût moyen est le plus faible pour
$\res{20 \text{ lampes}}$ par jour ($40$~€ par lampe).}{$c_{19}\approx 40{,}05$ et
$c_{21}\approx 40{,}05$}
\cas{4.}
\explique{On traduit « $c_{n}>50$ » par une inéquation du second degré :}{$n\ge 1$, donc
$n>0$}
\begin{calculs}
n+\dfrac{400}{n} &> 50 &&\\
n^{2}+400 &> 50n &&\justif{on multiplie par $n>0$ : le sens ne change pas}\\
n^{2}-50n+400 &> 0 &&\justif{on retire $50n$}
\end{calculs}
\begin{calculs}
\Delta &= (-50)^{2}-4\times 1\times 400 &&\justif{$a=1$, $b=-50$, $c=400$}\\
\Delta &= 2\,500-1\,600 = 900 &&\justif{$\sqrt{900}=30$}\\
x_{1} &= \dfrac{50-30}{2} = 10 &&\\
x_{2} &= \dfrac{50+30}{2} = 40 &&
\end{calculs}
\explique{$a=1>0$ : le trinôme est positif à l'extérieur des racines, soit $n<10$ ou
$n>40$.}{signe de $a$ à l'extérieur des racines}
\explique{Pour $n\ge 20$, il reste $n>40$, c'est-à-dire $n\ge 41$.}{$n$ est entier}
\explique{Le coût moyen dépasse $50$~€ à partir de $\res{41 \text{ lampes}}$ par
jour.}{$c_{40}=40+10=50$ ne dépasse pas ; $c_{41}\approx 50{,}76$}
\end{exo}

\end{document}
